\documentclass[11pt]{article}

\usepackage[margin=1in]{geometry}
\usepackage{amsmath,amssymb,amsthm,mathtools}
\usepackage{booktabs,tabularx,array}
\usepackage{graphicx}
\usepackage{microtype}
\usepackage{enumitem}
\usepackage{xcolor}
\usepackage{xurl}
\usepackage[hidelinks]{hyperref}

\definecolor{kernelcolor}{RGB}{0,92,74}
\definecolor{humanproofcolor}{RGB}{45,76,135}
\definecolor{opencolor}{RGB}{145,76,0}
\definecolor{classicalcolor}{RGB}{75,75,75}
\newcommand{\Kernel}{\textcolor{kernelcolor}{\textsc{Kernel}}}
\newcommand{\Human}{\textcolor{humanproofcolor}{\textsc{Proof}}}
\newcommand{\Open}{\textcolor{opencolor}{\textsc{Open}}}
\newcommand{\Classical}{\textcolor{classicalcolor}{\textsc{Classical}}}
\newcommand{\NewResult}{\textsc{New}}
\newcommand{\Realization}{\textsc{Realization}}
\newcommand{\NoGo}{\textsc{No-Go}}

\newcommand{\ii}{\mathrm{i}}
\newcommand{\tr}{\operatorname{tr}}
\newcommand{\rank}{\operatorname{rank}}
\newcommand{\diag}{\operatorname{diag}}
\newcommand{\C}{\mathbb{C}}
\newcommand{\R}{\mathbb{R}}
\newcommand{\id}{\mathbf{1}}
\newcommand{\norm}[1]{\left\lVert #1\right\rVert}

\newtheorem{theorem}{Theorem}[section]
\newtheorem{proposition}[theorem]{Proposition}
\newtheorem{corollary}[theorem]{Corollary}
\newtheorem{definition}[theorem]{Definition}
\newtheorem{remark}[theorem]{Remark}
\newtheorem{conjecture}[theorem]{Conjecture}

\graphicspath{{Sources/figures/}{figures/}}

\title{Null-Spinor Area as the Rest Gap of an Exactly Unitary Dirac Walk\\
\large A Machine-Verified Finite Construction}
% RELEASE GATE: replace the project author line with named authors,
% affiliations, ORCIDs, and a contribution statement before circulation.
\author{Mark Schwab}
\date{July 10, 2026}

\begin{document}
\maketitle

\begin{abstract}
The rest gap is normally an independent input to a Dirac quantum walk.  Here
it is fixed by supplied null-spinor geometry, with no second mass parameter,
and we prove exactly what that derivation adds.  For
massless rank-one momenta $P=\sum_i\psi_i\psi_i^\dagger$, we prove
\[
  \mu^2:=\det P
  =\sum_{i<j}|\psi_i\wedge\psi_j|^2,
\]
with $\mu=0$ exactly when the occupied spinors span at most one complex
direction.  For two
spinors, the complex area $z=\psi_1\wedge\psi_2$ determines the odd Hermitian
operator
\[
  B_z=\begin{pmatrix}0&z\\ \overline z&0\end{pmatrix},
  \qquad B_z^2=\det(P)\id,
  \qquad (k\sigma_z+B_z)^2=(k^2+\det P)\id,
\]
and every odd Hermitian rest operator is $B_w$ for a unique $w\in\C$: the rest
gap is the null-spinor area, with no independent mass parameter.
Exponentiating $B_z$ gives an exactly unitary walk with quasienergy
$\cos(\omega a)=\cos(ka)\cos(\mu a)$, corner ratio $-\ii\tan(a\mu)$, and a
directed history sum in which opposite turns carry $z$ and $\bar z$.  We prove
that the walk's fixed-momentum symbols converge to the corresponding Dirac
flow uniformly on bounded momentum boxes at rate $O(1/n)$, for the $1+1$
walk and for an ordered $3+1$ successive-axis walk with the complex phase
retained; an arbitrary position-dependent Pluecker field $z(x)$ generates an
exactly reversible finite-torus walk with a strict causal cone, site-dependent
chiral rotation induces explicit endpoint links while a constant phase induces
none, and compactly supported momentum-space data converge in $L^2$ at every
fixed time.  The link data derived from any nowhere-zero Pluecker field has
an exact $2\pi$-integer total turning, invariant under global chiral phase,
with winding one realized by an explicit field generated by primitive
spinors; the patched data is derived rather than supplied, and the
companion defect analysis gives an exact involutive-compression mode mechanism
and explicit mode/no-mode fixtures, but a complete four-site sign-family audit
proves that winding alone does not determine the compression signature: wall
position relative to the reflection-fixed sites is also required.  A constant
$|z|$ alone would only reparametrize an assigned Dirac
mass; the construction escapes through its oriented data.  The gap-closing
locus is the geometric collinearity condition, the directed weights and the
local connection distinguish $z$ from $|z|$, and a finite quartic pair
transfer, together with its exactly unitary Fock-norm-preserving pair kick,
has forward two-particle amplitude $z$ and reverse amplitude $\bar z$.  The
same kick admits an exact placement on any embedded four-mode set as an even
quartic CAR polynomial.  Disjoint sets commute, and a finite unit-phase gate
schedule preserves strong CAR support inside the union of its declared sets.
A graph-metric successor proves more: for a sequential list of finite-range
blocks, support remains inside the corresponding iterated neighborhood, with
an explicit contiguous witness and a nonlocal control.  A further exact
layer theorem shows that pairwise-disjoint local gates cost only one
neighborhood expansion per circuit layer, so an arbitrary finite schedule is
bounded by its layer depth rather than its gate count.  The
constant phase is invisible to the one-particle rest spectrum --- the rest
operators of $3+4\ii$ and $5$ are exactly unitarily conjugate --- but the
\emph{transported relative} phase is already a free-carrier observable: for
two sites of equal modulus coupled by a chirality-sensitive kinetic edge,
$(H_\chi^2-(m^2+t^2)\id)^2=t^2\,|z_R-e^{-\ii\chi}z_L|^2\,\id$ exactly, so
equal-modulus fields with different transported phase have different
spectra, with the equal-modulus hypothesis load-bearing (an explicit
unequal-modulus counterexample refutes the identity without it), an exact
zero-mode locus $t=|z_L|$, $z_R=-e^{-\ii\chi}z_L$, and exact removability of
any common phase.  The phase also appears operationally in a supplied
two-particle pair-kick sector, with return probability $4/5$ versus $1$ at
those witnesses; the supplied kick is itself structured --- it equals, up
to a global phase, the exact quarter-period pulse of a finite Hermitian
even quartic CAR generator whose cube closes ($K^3=|z|^2K$), so the
one-parameter gate family is exactly integrable --- but whether the free
walk forces that interaction remains open, and the generator is supplied,
not derived.  The
same framework proves its own boundaries.  An exact all-zone determinant
factorization, $\det(U^{(4)}\mp\id)=4P_{0,\pi}$ with explicit polynomials,
holds at every momentum and mass angle, and for every massive angle with
$0<|\cos\theta|<1$ it yields a complete crossing classification: zero- and
$\pi$-quasienergy modes occur exactly where all three momentum cosines
vanish simultaneously.  The massless cubic regulator has
exactly three non-origin zero-quasienergy corner aliases; the massive
body-center step retains $\pm1$ eigenmodes for every mass angle, so this
regulator cannot open a global Floquet gap; no degree-one nearest-neighbor
factor supports a stationary amplitude compatible with origin normalization,
exact unitarity, and the full Dirac tangent; every momentum-independent onsite
coin leaves all three even-parity aliases intact; the positive homogeneous
Pluecker action selects only the collapsed zero scale; and exact two-step
blocking closes the mass subgroup but not the ordered split family.  A
pre-registered high-momentum benchmark passes all alias, body-center,
Wilson-gap, and negative-control thresholds.  Every finite statement above is
checked in Lean~4; the novelty is the derived rest operator and its verified
dynamical chain, not a new constant-coefficient $1+1$ walk.  A strict
alias-free discrete-time $3+1$ regulator, position-space PDE convergence, a
local interacting Fock dynamics, and any absolute mass scale remain open.
\end{abstract}

\paragraph{How to read the status marks.}
Origin and verification are independent.  \Classical{} identifies established
mathematics formalized here; \NewResult{} identifies a new theorem or
construction; \Realization{} identifies an explicit model.  \Kernel{} means
the stated finite claim is checked by the pinned Lean kernel; \Human{} means a
complete mathematical proof is given here but its dedicated Lean declaration
is pending; \Open{} marks a named unsolved problem.  Physical interpretations
are never promoted by a verification mark.

\section{Introduction: one invariant, one operator, one walk}

The sum of lightlike momenta need not be lightlike.  In the spinor
representation, each future null momentum is a rank-one positive matrix
$\psi\psi^\dagger$.  A collection of such momenta remains on the null boundary
exactly while its spinors occupy one complex direction.  Once two directions
open, their exterior area is nonzero and the sum moves into the timelike
interior.

This familiar geometry contains a sharper operator statement.  The complex
area $z=\psi\wedge\phi$ is not merely a number whose squared modulus happens
to equal mass squared.  It is the off-diagonal entry of a canonical odd
Hermitian operator.  That operator squares to the invariant momentum
determinant, anticommutes with the velocity grading, and therefore supplies the
rest term of a finite Dirac Hamiltonian.  Its exponential is an exact unitary
turn gate.  The same object thus controls quadratic cost, first-order rest
energy, finite-time rotation, and corner amplitudes.

\begin{figure}[t]
\centering
\includegraphics[width=0.98\linewidth]{null_spinor_architecture.pdf}
\caption{The complete finite construction.  The input is a pair of null
spinors, not a preassigned mass angle.  The complex area fixes both the rest
spectrum and the orientation phase of the unitary turn; collinearity is the
exact gap-closing boundary.}
\label{fig:pipeline}
\end{figure}

\begin{center}
\fbox{\begin{minipage}{0.92\linewidth}
\textbf{Synthesis theorem.}
Let $\psi,\phi\in\C^2$, set
$P=\psi\psi^\dagger+\phi\phi^\dagger$,
$z=\psi\wedge\phi$, and $\mu=|z|$.  Then
\[
 \det P=\mu^2,
 \qquad
 B_z=\begin{pmatrix}0&z\\ \overline z&0\end{pmatrix},
 \qquad
 B_z^2=\mu^2\id.
\]
Writing $\Gamma=\diag(1,-1)$ and $H_z(k)=k\Gamma+B_z$, one has
\[
 \Gamma B_z+B_z\Gamma=0,
 \qquad
 \boxed{H_z(k)^2=(k^2+\det P)\id}.
\]
Thus the positive rest energy is $\mu=\sqrt{\det P}$, and the gap closes
exactly on the collinear null boundary.  Multiplying $z$ by a unit phase
changes $B_z$ by unitary conjugation, so the spectrum depends only on $P$ while
the oriented decomposition retains a path phase.  \NewResult{} \Kernel{}
\end{minipage}}
\end{center}

\begin{center}
\fbox{\begin{minipage}{0.92\linewidth}
\textbf{What changes physically, and what does not.}
The Pluecker origin fixes more than a dispersion parameter: it fixes the
complex directed turn amplitudes, their local chiral connection under a
spatially varying field, the exact gap-closing locus, and a finite quartic
pair-transfer amplitude that distinguishes $z$ from $\bar z$.  It does not by
itself cure lattice aliases or predict an absolute mass.  In fact, within the
four-channel separable degree-one architecture, every even-parity zone corner
is provably locked to the origin after any onsite Pluecker coin; under
the natural positive homogeneous action, every positive configuration has
only the collapsed stationary scale.  These are falsifiable structural
consequences, not interpretive qualifications.  \NewResult{} \Kernel{}
\end{minipage}}
\end{center}

The contributions form two pillars.  The first is a derivation: the
canonical rest operator and its verified dynamical chain, from spinor pair
to walk to histories to uniform continuum control.  The second is an
obstruction suite: exact all-zone criteria and no-go theorems that fix the
boundary of the minimal architecture --- including a corollary constraining
the newest published doubler-free constructions --- and that turn the search
for a clean strict $3+1$ regulator into a precise algebra problem.  Seven
contributions organize the paper.
\begin{enumerate}[leftmargin=2em]
  \item \textbf{The invariant.}  We isolate the exterior-area invariant, its
  exact massless boundary, and the grading between the first-order gap $\mu$
  and the quadratic invariant $\mu^2$.  The Cauchy--Binet identity is
  classical; its complete finite formalization and controls are part of the
  contribution.
  \item \textbf{The canonical operator.}  We construct $B_z$ directly from the
  complex Pluecker coordinate, prove the synthesis theorem for arbitrary
  spinor pairs, and classify every odd Hermitian rest operator as $B_w$ for a
  unique $w$.  No rational example or separately chosen Dirac mass parameter
  carries the general statement.
  \item \textbf{Exact histories.}  We derive the exact unitary turn gate and
  its corner expansion, with finite-spacing corner parameter $\tan(a\mu)$
  rather than $a\mu$.  A directed-edge path theorem places $z$ and $\bar z$ in
  the same finite history sum and proves equality with the actual complex
  Pluecker transfer power.
  \item \textbf{Uniform continuum control.}  We prove kernel-checked, uniform
  bounded-momentum $O(1/n)$ many-step estimates for the $1+1$ and ordered
  $3+1$ split-step symbols.  For the explicit schedule $n=K^4$, the $3+1$
  estimate remains convergent on arbitrary parameter sequences satisfying
  $B_4\leq K$, with an $O(K^{-2})\exp(|t|/K^3)$ envelope.  The full complex phase is covered by unitary
  conjugacy and by the direct complex-coin composition, not discarded after
  taking its modulus.
  \item \textbf{A finite local $3+1$ walk with exact Fourier control.}  We
  construct an exactly norm-preserving finite-torus walk with an explicit
  two-sided inverse and a strict one-shift-per-axis causal cone, whose ordered
  symbol is the product controlled by the $3+1$ estimate.  The
  finite-to-analytic momentum-sign bridge is exact, an iterated three-torus
  Plancherel theorem gives the exact $N^{-3}$ wave-packet normalization, and a
  compact-support multiplier theorem yields momentum-space $L^2$ convergence
  at every fixed time.  Arbitrary position-dependent Pluecker fields evolve
  unitarily, and site-dependent chiral rotation induces exact endpoint links.
  \item \textbf{The obstruction audit.}  The same framework proves where it
  stops, and now does so with an exact all-zone criterion: the determinant
  factorizations $\det(U^{(4)}\mp\id)=4P_{0,\pi}$ hold at every real momentum
  and mass angle, so zero and $\pi$ modes are exactly the zeros of two
  explicit polynomials --- completely classified for $0<|\cos\theta|<1$ as
  the simultaneous vanishing of all three momentum cosines.  The massless
  cubic regulator has three nonzero
  zero-quasienergy
  corner aliases; the massive body-center step retains exact $+1$ and $-1$
  modes for every mass angle, so the regulator is globally ungapped; a
  degree-one nearest-neighbor factor admits no stationary amplitude compatible
  with origin normalization, exact unitarity, and the full Dirac tangent; and
  every momentum-independent onsite coin leaves all three even-parity aliases
  intact.  Exact temporal blocking closes the mass subgroup but generates new
  ordered words in the full split family, and a homogeneous-action no-go
  proves that selecting a nonzero $|z|$ requires an additional scale-bearing
  input.
  \item \textbf{The many-body seed.}  A generic occupation-basis construction
  satisfies all canonical anticommutation relations, and a four-mode quartic
  pair transfer has oriented two-particle amplitudes $z$ and $\bar z$, exactly
  distinguished at the fixture $z=3+4\ii$.  The associated normalized pair
  kick is an explicit complex-linear automorphism, involutive and
  Fock-inner-product preserving; the Hermitian quartic operator equals the
  kick on the pair sector and the identity off it, and the kick fixes every
  one-particle basis state while acting nontrivially on the displayed
  two-particle state.  An interference corollary shows that the supplied
  interaction sector retains the phase:
  equal-modulus fields with exactly conjugate one-particle rest operators
  are separated, inside the supplied pair-kick sector, by a double-kick
  return probability of $4/5$ versus $1$.  This is not yet a second quantization of the spatial
  walk: the determinant-minor lift now has exact vacuum, one-particle, number,
  parity, creation-covariance, identity, and composition theorems, including
  $\Gamma(UV)=\Gamma(U)\Gamma(V)$, and its ordered-minor coefficients obey
  exact conjugate-transpose compatibility.  A one-particle unitary's lifted
  conjugate transpose is an exact two-sided inverse, and the lift preserves
  the finite occupation-basis Hermitian inner product.  Creation and
  annihilation covariance are exact.  Separately, the pair kick can be placed
  on arbitrary four-mode sets; disjoint gates commute, and every finite
  declared schedule has exact support propagation in the even CAR algebra,
  with a nonzero phase-transfer witness inside the acted set.  Under an explicit
  graph neighborhood and finite-range condition on every block, a sequential
  schedule obeys an exact iterated-neighborhood cone.  Pairwise-disjoint local
  gates have now been grouped into layers, with one neighborhood expansion per
  layer and an exact radius bound by circuit depth.  Composing this
  interaction layer with the spatial free walk, deriving an
  operator-exponential law, and computing scattering or binding remain open.
\end{enumerate}

The result sits between spinor-helicity kinematics and discrete relativistic
quantum dynamics \cite{ElvangHuang,ArkaniHamedMassive}.  Cauchy--Binet and the
Pauli map supply the classical foundation.  Feynman's checkerboard and modern
quantum walks show how a massive Dirac propagator can arise from null steps and
direction changes
\cite{FeynmanHibbs,JacobsonSchulman,BisioDarianoTosini,DarianoPathIntegral,MlodinowBrun}.
Our contribution is the exact finite map from the oriented spinor area to the
odd rest operator and then to a controlled unitary walk, together with a
machine-checked dependency boundary.

\subsection*{Nearest constructions and the exact novelty claim}

For constant $z$, the update in this paper is unitarily equivalent to the
standard $1+1$ split-step Dirac walk.  We do not claim otherwise.  What is new
is the input map and its naturality: the coin is generated by the complex
Pluecker coordinate of null momentum factors, its modulus gives the rest
energy, its phase acts by internal conjugacy and survives in oriented turn
amplitudes, and the resulting dependency chain is machine checked.  Table
\ref{tab:comparison} positions that claim against the closest constructions.

The Pavia programme moreover derives the Weyl and Dirac automata from
informational first principles \cite{FreeQCA}; the derivation offered here
is geometric rather than axiomatic, and the two are answers to different
questions.

The kinematic layer itself has long roots, and we claim none of it.
Weyl--van der Waerden and massive spinor-helicity formalisms routinely
decompose a massive momentum into two null spinors, with
$\det K=\lambda_1\lambda_2=k^2$ appearing explicitly in the
eigen-decomposition of the momentum bispinor \cite{DittmaierWvdW}, and the
bold-spinor formalism of Arkani-Hamed, Huang, and Huang
\cite{ArkaniHamedMassive} makes the two-null-spinor description of massive
states the organizing principle of amplitude kinematics.  Those formalisms
use the decomposition as covariant bookkeeping for scattering amplitudes.
The present paper starts where that bookkeeping stops: the wedge coordinate
is promoted to the off-diagonal entry of a canonical odd Hermitian operator,
the operator is classified (Theorem~\ref{thm:classification}), exponentiated
into an exactly unitary walk, tracked through directed histories and a local
connection, and every finite statement in that chain is kernel-checked.  The
amplitude formalisms neither assert nor need such a dynamical chain; the
overlap is the classical $2\times2$ Cauchy--Binet layer, which we label
\Classical{} throughout.

\begin{table}[t]
\centering
\scriptsize
\resizebox{\linewidth}{!}{%
\begin{tabular}{@{}lcccccc@{}}
\toprule
Construction & Exact local unitary & Exact finite histories & Mass input
  & Dimension & Continuum control & Machine checked\\
\midrule
Feynman checkerboard \cite{FeynmanHibbs,JacobsonSchulman}
  & not its object & yes & assigned corner weight & $1+1$ & propagator limit & no\\
Dirac QCA \cite{BisioDarianoTosini,DarianoPathIntegral}
  & yes & yes & automaton parameter & $1+1$ & channel/wave-packet bounds & no\\
Split-step DCA \cite{MallickSplitStep}
  & yes & not central & tunable coin angle & $1+1$ & Dirac limit & no\\
Weyl/Dirac QCA \cite{FreeQCA}
  & yes & model dependent & Weyl coupling & up to $3+1$ & relativistic limit & no\\
Tetrahedral Dirac walk \cite{NzonganiTetra}
  & yes & not central & assigned $e^{-\ii\varepsilon m\gamma^0}$ coin
  & $3+1$ tetrahedra & first-order Dirac limit & no\\
Successive-axis walk \cite{MlodinowBrun}
  & yes & not central & assigned mass parameter
  & $3+1$ & first-order relativistic limit & no\\
This work
  & yes & yes & $z=\psi\wedge\phi$ & $1+1$; finite $3+1$ lift
  & uniform $O(1/n)$ symbols; three-torus $L^2$ & yes\\
\bottomrule
\end{tabular}}
\caption{Structural comparison, not a priority ranking.  The closest QCA
works have stronger field-theoretic or higher-dimensional scope.  The present
distinction is the geometrically generated complex rest operator, the exact
kernel dictionary, and their formal verification.}
\label{tab:comparison}
\end{table}

\subsection*{The reparametrization test}

At constant nonzero $z$, the one-particle spectrum alone cannot distinguish
this construction from a Dirac walk whose mass was assigned as $m=|z|$.  The
claim of added content must therefore survive after that spectrum is held
fixed.  Table~\ref{tab:reparametrization} states the test explicitly.

\begin{table}[h]
\centering
\small
\begin{tabularx}{\linewidth}{@{}p{0.22\linewidth}XX@{}}
\toprule
Question & Assigned scalar mass & Pluecker-derived input in this paper\\
\midrule
Constant dispersion
 & Same after setting $m=|z|$
 & No claimed distinction\\
Gap-closing locus
 & External condition $m=0$
 & Exact collinearity condition $\psi\wedge\phi=0$\\
Directed histories
 & A real or independently chosen turn coefficient
 & Opposite orientations carry $z$ and $\bar z$ in the same transfer sum\\
Spatially varying phase
 & Additional complex-mass field supplied independently
 & Primitive-spinor phase induces an exact endpoint connection under local chiral conjugacy\\
Finite two-particle probe
 & Requires an independent interaction coefficient
 & The exhibited quartic transfer reuses the same wedge as amplitudes $z$ and
 $\bar z$; equal-modulus fields with conjugate one-particle operators are
 separated by an exact interference probability ($4/5$ versus $1$) in the
 supplied pair-kick sector\\
Absolute scale
 & Not predicted
 & Not predicted; the homogeneous-action no-go proves an added scale is required\\
Ultraviolet lattice spectrum
 & Regulator dependent
 & Still regulator dependent; the present walk retains exact aliases\\
\bottomrule
\end{tabularx}
\caption{What is and is not gained by deriving the mass parameter.  The
quartic row is an explicit finite extension, not yet a theorem that the free
action uniquely forces that interaction.}
\label{tab:reparametrization}
\end{table}

This comparison also fixes the paper's ceiling.  Its centerpiece is a
one-particle $1+1$ walk with a new geometric input and a rigorous finite chain;
Section~\ref{sec:threeplusone} adds a controlled finite $3+1$ local extension.
It is not yet a new QCA classification, an interacting QCA, or a field theory;
 its exact high-symmetry audit in fact exposes corner aliases and a
 mass-independent body-center no-gap obstruction that a successor regulator
 must remove.

\section{Null-spinor area and the future cone}

\subsection{From rank-one momenta to area}

Let $\psi_1,\ldots,\psi_N\in\C^2$ and define
\begin{equation}\label{eq:momentum}
 M=\begin{bmatrix}\psi_1&\cdots&\psi_N\end{bmatrix},
 \qquad
 P=MM^\dagger=\sum_{i=1}^N\psi_i\psi_i^\dagger.
\end{equation}
Each summand is positive semidefinite and has rank at most one.  Under the
mostly-minus Pauli map, it is a future null four-momentum.  The determinant of
$P$ is the Lorentzian square of the total momentum.

\begin{theorem}[Null-spinor area identity]\label{thm:area}
For the matrix in \eqref{eq:momentum},
\begin{equation}\label{eq:area}
 \det P
 =\norm{\bigwedge{}^2 M}_{\mathrm{HS}}^2
 =\sum_{i<j}\left|\det\begin{bmatrix}\psi_i&\psi_j\end{bmatrix}\right|^2.
\end{equation}
Consequently $\det P\geq0$, and $\det P=0$ if and only if the occupied columns
span a subspace of dimension at most one.
\end{theorem}
\noindent\Classical{} \Kernel{}

\begin{proof}
Cauchy--Binet applied to $MM^\dagger$ expresses its determinant as the sum of
the squared moduli of all $2\times2$ minors of $M$.  Those minors are exactly
the coordinates of $\bigwedge^2M$.  Every summand is nonnegative, so the sum
vanishes exactly when every pairwise wedge vanishes, equivalently when
$\rank M\leq1$.
\end{proof}

Define
\begin{equation}\label{eq:mu}
 \mu:=\norm{\bigwedge{}^2M}_{\mathrm{HS}}=\sqrt{\det P}.
\end{equation}
The quantity $\mu$ is the first-order rest energy; $\mu^2$ is the Lorentzian
quadratic invariant.  The exterior norm makes the geometric word ``area''
literal rather than metaphorical.

For readers who prefer spatial directions, write the null momenta as
$p_i=E_i(1,\mathbf n_i)$ with $E_i\geq0$ and $|\mathbf n_i|=1$.  With the same
mostly-minus normalization,
\begin{equation}\label{eq:angular}
 \mu^2
 =2\sum_{i<j}E_iE_j\bigl(1-\mathbf n_i\mathbin{\cdot}\mathbf n_j\bigr).
\end{equation}
Mass squared is therefore a positive sum of directional disagreements.  It is
zero only when every occupied direction agrees; no cancellation between
different pairs is available.

\subsection{Why the determinant is canonical}

The determinant is not selected merely because it gives the desired answer.
On the physical space of Hermitian $2\times2$ matrices it is the unique real
quadratic form, up to scale, that vanishes on the full rank-one positive
boundary.

\begin{proposition}[Null-boundary uniqueness]\label{prop:unique}
Let $q$ be a real quadratic form on $\operatorname{Herm}(2)$.  If
$q(\psi\psi^\dagger)=0$ for every $\psi\in\C^2$, then
$q(P)=c\det P$ for some $c\in\R$.
\end{proposition}
\noindent\Classical{} \Kernel{}

\begin{proof}
Identify $P\in\operatorname{Herm}(2)$ with $(t,\mathbf x)\in\R^{1,3}$ so that
$\det P=t^2-|\mathbf x|^2$.  Write
$q(t,\mathbf x)=a t^2+2t\,\mathbf b\cdot\mathbf x+\mathbf x^T C\mathbf x$,
with $C$ symmetric.  On the future null cone, homogeneity reduces the
hypothesis to $q(1,\mathbf n)=0$ for every unit vector $\mathbf n$.  Comparing
$\mathbf n$ with $-\mathbf n$ gives $\mathbf b=0$.  Hence
$\mathbf n^TC\mathbf n=-a$ on the unit sphere, which forces $C=-aI$ by
polarization.  Thus $q=a(t^2-|\mathbf x|^2)=a\det P$.
\end{proof}

The Lean theorem uses the four real coordinates of a Hermitian matrix and nine
explicit rank-one spinor probes to eliminate every quadratic coefficient except
one multiple of $pq-r^2-s^2$.  The older real-symmetric slice and a form that
fails the full-boundary hypothesis remain as independent controls.

\begin{corollary}[First-order uniqueness]
Suppose $g$ is nonnegative and homogeneous of degree one on the future cone,
and $g(P)^2$ is a quadratic form vanishing on the rank-one boundary.  Then
$g(P)=c\sqrt{\det P}$ for some $c\geq0$.
\end{corollary}
\noindent\Classical{} \Human{}

This is the logical bridge from a canonical mass-squared invariant to a
canonical rest gap.  The normalization $c=1$ fixes units.

\section{The Pluecker mass operator}

For two spinors, retain the complex orientation data rather than immediately
discarding it into an absolute square:
\begin{equation}
 z:=\psi\wedge\phi=\psi_0\phi_1-\psi_1\phi_0,
 \qquad
 P:=\psi\psi^\dagger+\phi\phi^\dagger.
\end{equation}
Theorem~\ref{thm:area} gives $|z|^2=\det P$.  Define
\begin{equation}\label{eq:Bz}
 B_z:=\begin{pmatrix}0&z\\ \overline z&0\end{pmatrix},
 \qquad
 \Gamma:=\begin{pmatrix}1&0\\0&-1\end{pmatrix}.
\end{equation}

\begin{theorem}[Pluecker-derived rest operator]\label{thm:Bz}
The matrix $B_z$ is Hermitian and odd for the grading $\Gamma$.  More
precisely,
\begin{equation}\label{eq:Bzidentities}
 B_z^\dagger=B_z,
 \qquad
 \Gamma B_z+B_z\Gamma=0,
 \qquad
 B_z^2=|z|^2\id=\det(P)\id.
\end{equation}
For
\begin{equation}\label{eq:Hz}
 H_z(k):=k\Gamma+B_z,
\end{equation}
one has
\begin{equation}\label{eq:hero}
 \boxed{H_z(k)^2=\left(k^2+\det P\right)\id}.
\end{equation}
The rest eigenvalues are $\pm|z|$, and $B_z=0$ exactly when the two nonzero
spinor directions are collinear.
\end{theorem}
\noindent\NewResult{} \Kernel{}

\begin{proof}
Hermiticity is immediate from the two off-diagonal entries.  Direct
multiplication gives
$B_z^2=\diag(z\overline z,\overline z z)=|z|^2\id$, while the diagonal signs
in $\Gamma$ give the anticommutator zero.  Expanding $H_z(k)^2$ leaves
$k^2\id+B_z^2$ because the mixed term vanishes.  Finally $B_z=0$ if and only
if $z=0$, and a two-spinor wedge vanishes exactly on projective collinearity.
Every step of this proof is represented by a general Lean declaration, not by
evaluation of a fixed example.
\end{proof}

Theorem~\ref{thm:Bz} is stronger than assigning the scalar $\det P$ to an
otherwise independent Dirac mass slot.  The odd operator is constructed from
the oriented Pluecker coordinate itself; its square derives the quadratic
invariant.  The distinction between powers is now fixed:

\begin{center}
\begin{tabularx}{0.92\linewidth}{@{}X>{\raggedleft\arraybackslash}p{0.38\linewidth}@{}}
\toprule
Layer & Natural quantity\\
\midrule
Null-spinor kinematics & $\mu^2=\det P=|z|^2$\\
Odd mass operator & $B_z$, with $B_z^2=\mu^2\id$\\
Quadratic costs and Hessians & $\mu^2$\\
Dirac rest energy & $\mu=|z|$\\
Finite unitary turn angle & $a\mu$\\
Exact corner/stay ratio & $\tan(a\mu)$\\
\bottomrule
\end{tabularx}
\end{center}

\subsection{Phase covariance and decomposition independence}

If $u\in\C$ has $|u|=1$, let $W_u=\diag(u,1)$.  Then
\begin{equation}\label{eq:phasecov}
 W_uB_zW_u^\dagger=B_{uz},
 \qquad
 W_uW_u^\dagger=\id.
\end{equation}
Both identities are kernel-checked.  Thus the phase of $z$ rotates the internal
basis but cannot change the rest spectrum.

The construction is also covariant under the spacetime spin group.  If
$A\in\mathrm{SL}(2,\C)$ acts on both spinors by left multiplication, then
\begin{equation}\label{eq:sl2cov}
 (A\psi)\wedge(A\phi)=\det(A)(\psi\wedge\phi)=z,
 \qquad
 B_{(A\psi)\wedge(A\phi)}=B_z.
\end{equation}
This is the algebraic $\mathrm{SL}(2,\C)$ statement covering
proper-orthochronous Lorentz covariance; it is \Kernel{}.  General
$A\in\mathrm{GL}(2,\C)$ acts with the expected determinant weight; that
weighted extension is an algebraic \Human{} corollary, not a separate Lean
declaration in the present development.

\begin{theorem}[Classification of odd Hermitian rest operators]
\label{thm:classification}
Let $B\in M_2(\C)$ be Hermitian and satisfy
$\Gamma B+B\Gamma=0$.  Then there is a unique $w\in\C$ such that
\begin{equation}
 B=B_w=\begin{pmatrix}0&w\\\overline w&0\end{pmatrix},
 \qquad
 B^2=|w|^2\id.
\end{equation}
Consequently the odd Hermitian square roots of $\mu^2\id$ form one phase
circle $w=\mu u$, $|u|=1$, and for $\mu>0$ they are all related by the
conjugacies in \eqref{eq:phasecov}.  Specifying the oriented Pluecker
coordinate $z$ selects one member of that orbit.
\end{theorem}
\noindent\NewResult{} \Kernel{} for the normal form and square;
\Human{} for the phase-orbit corollary.

\begin{proof}
Oddness forces both diagonal entries of $B$ to vanish, and Hermiticity forces
the lower-left entry to be the complex conjugate of the upper-right entry
$w=B_{01}$.  The square is then $|w|^2\id$.  If it equals $\mu^2\id$, then
$|w|=\mu$; for $\mu>0$, write $w=\mu u$ with $|u|=1$ and apply
\eqref{eq:phasecov}.
\end{proof}

This also controls the nonuniqueness of the null decomposition.  If
$M=[\psi\ \phi]$ has full rank and $M'=MR$ with $R\in U(2)$, then
$M'M'^\dagger=MM^\dagger$ and
$\det M'=(\det R)\det M$.  Since $|\det R|=1$, the two mass operators are
related by \eqref{eq:phasecov}.  The unitary-equivalence class of $B_z$ depends
only on $P$; an oriented factorization retains the phase needed by individual
turn amplitudes.  \Classical{} \Kernel{}

The distinction matters physically.  The spectrum is decomposition
independent, while an amplitude can remain sensitive to orientation.  The
construction preserves exactly the information needed by a path sum and
forgets exactly the information a rest energy cannot observe.

\section{Exact unitary null histories}

\subsection{Exponentiating the derived rest term}

For $\mu=|z|>0$, set $\widehat B_z=B_z/\mu$.  Since
$\widehat B_z^2=\id$,
\begin{equation}\label{eq:expB}
 e^{-\ii aB_z}
 =\cos(a\mu)\id-\ii\sin(a\mu)\widehat B_z.
\end{equation}
At $\mu=0$ the expression is understood by continuity and equals $\id$.
Combining the turn with the conditional null translation gives the momentum
symbol
\begin{equation}\label{eq:walk}
 U_a(k;z):=e^{-\ii ak\Gamma}e^{-\ii aB_z}.
\end{equation}
Each factor is unitary.  The two components translate in opposite null
directions; only the odd factor changes direction.

The finite formula in \eqref{eq:expB} is itself closed under composition:
writing it as $C_z(a)$, one has
\begin{equation}
 C_z(a)C_z(a)^\dagger=\id,
 \qquad
 C_z(a)C_z(b)=C_z(a+b).
\end{equation}
The explicit rest eigenvectors, right-unitary decomposition theorem, and this
one-parameter group law are all \Kernel{}.

By phase covariance, $U_a(k;z)$ is unitarily conjugate to the real-mass walk
with parameter $\mu$.  Its determinant is one and its trace is
$2\cos(ka)\cos(\mu a)$.  Therefore an eigenphase $e^{-\ii\omega a}$ obeys
\begin{equation}\label{eq:dispersion}
 \boxed{\cos(\omega a)=\cos(ka)\cos(\mu a)}.
\end{equation}
At $k=0$ the exact eigenphases are $e^{\mp\ii a\mu}$, so the positive
principal-branch quasienergy is $\mu$.  Quasienergy is periodic modulo
$2\pi/a$; no global unbranched equality is asserted.  The real-mass spectral
statement and oriented mass-coin group law are \Kernel{}; the displayed
spectral conjugacy is an immediate corollary of the kernel-checked phase
covariance.

\subsection{The exact corner ratio}

Equation~\eqref{eq:expB} gives two kinds of matrix element.  A component that
keeps its direction contributes $\cos(a\mu)$.  A direction change contributes
$-\ii\sin(a\mu)$ times $z/\mu$ or $\overline z/\mu$, according to its
orientation.  Hence an $n$-step history $h$ with $c(h)$ turns has amplitude
\begin{equation}\label{eq:pathweight}
 A_a(h)
 =e^{\ii\Phi_k(h)}
   \cos(a\mu)^{n-c(h)}
   \bigl(-\ii\sin(a\mu)\bigr)^{c(h)}
   \Xi_z(h),
\end{equation}
where $\Phi_k(h)$ is the translation phase and $\Xi_z(h)$ is a product of
unit-modulus orientation factors.  Whenever $\cos(a\mu)\neq0$,
\begin{equation}\label{eq:tanweight}
 A_a(h)
 =e^{\ii\Phi_k(h)}\cos(a\mu)^n
   \bigl(-\ii\tan(a\mu)\bigr)^{c(h)}\Xi_z(h).
\end{equation}

\begin{proposition}[Unitary checkerboard correspondence]\label{prop:tan}
After removal of the common factor $\cos(a\mu)^n$, the exactly unitary walk is
a corner-weighted checkerboard sum with finite-spacing corner parameter
$\tan(a\mu)$ and the orientation phase inherited from $z$.  Under the opposite
checkerboard sign convention, the parameter changes sign.  In either
convention,
\begin{equation}
 \tan(a\mu)=a\mu+O(a^3),
\end{equation}
but $\tan(a\mu)$ and $a\mu$ are not identical at finite spacing.
\end{proposition}
\noindent\Classical{} for the coin parametrization; \NewResult{} \Kernel{} for
the exact two-weight recursion scaling and the separate transfer-matrix
identification; \NewResult{} \Kernel{} for the directed complex orientation
bookkeeping in the same path sum.

\begin{proof}
Insert the two matrix elements from \eqref{eq:expB} at every time slice.  A
history with $c(h)$ off-diagonal selections and $n-c(h)$ diagonal selections
has exactly the product in \eqref{eq:pathweight}.  Factoring one
$\cos(a\mu)$ from every slice gives \eqref{eq:tanweight}.  The Taylor statement
is the standard expansion of the tangent.
\end{proof}

The Lean development verifies the complete finite dictionary.  Its recursive
two-weight kernel satisfies
\[
 K[s,sw]_n=s^nK[1,w]_n,
\]
and specialization to $s=\cos(a\mu)$ and
$w=-\ii\tan(a\mu)$ gives the unitary checkerboard correspondence exactly when
$\cos(a\mu)\neq0$.  A quarter-angle nondegenerate witness and the
$a\mu=\pi/2$ pure-flip boundary are included.  Separately, the development
verifies that finite null histories weighted by a corner power equal a closed
binomial kernel and satisfy a discrete Dirac recursion, and that the same
direction-resolved history sum equals the corresponding transfer-matrix power.
The direct Gaussian-rational path-sum declaration uses a real mass and corner
weight $+\ii\varepsilon m$ in its stated convention.  Independently, the
directed-edge theorem assigns every history edge the corresponding entry of
the actual complex Pluecker mass coin and proves that the complete oriented
history sum equals its transfer-matrix power.  Thus $\Xi_z$ in
\eqref{eq:pathweight} is now packaged directly: opposite turns carry $z$ and
$\bar z$, with the exact $3+4\ii$ control proving that the two normalized
orientation factors are nonzero and unequal.

The tangent appearance is not asserted as novel by itself.  Exponential and
Euler-angle parametrizations of unitary Dirac-walk coins were related earlier,
including explicit tangent formulas in quantum lattice-Boltzmann work
\cite{SucciQW}.  The contribution here is the exact recursive-kernel scaling
law composed with a coin whose generator is derived from the complex Pluecker
coordinate, together with its machine verification.

\section{Full lattice spectrum, symmetry, and doubling audit}
\label{sec:spectrum}

Quantum-walk claims cannot be judged only at small momentum.  Introduce the
dimensionless variables
\[
 q=ka\in[-\pi,\pi],
 \qquad
 \theta=\mu a,
 \qquad
 \Omega=\omega a.
\]
The Brillouin zone is periodic in $q$, and quasienergy is periodic modulo
$2\pi$.

\begin{theorem}[Complete two-band audit]\label{thm:bands}
The two quasienergy bands are
\begin{equation}\label{eq:bands}
 \Omega_\pm(q)
 =\pm\arccos\!\bigl(\cos q\cos\theta\bigr)
 \pmod{2\pi}.
\end{equation}
They are degenerate only when $q\in\pi\mathbb Z$ and
$\theta\in\pi\mathbb Z$.  On the physical principal branch
$0<|\theta|<\pi/2$, the walk is gapped throughout the Brillouin zone: the
positive rest energy at $q=0$ is $|\theta|/a=\mu$, the interband zero-energy
gap is $2\mu$, and the circular quasienergy gap at the zone edge is also
$2|\theta|$.  At $\theta=0$ there are two massless cones:
\begin{equation}
 (q,\Omega)=(0,0)
 \qquad\text{and}\qquad
 (q,\Omega)=(\pi,\pi).
\end{equation}
For a differentiable band away from a degeneracy,
\begin{equation}\label{eq:groupvelocity}
 v_g=\frac{\partial\omega}{\partial k}
 =\frac{\sin q\cos\theta}{\sin\Omega},
 \qquad |v_g|\leq1,
\end{equation}
with strict inequality for a massive angle whenever the denominator is
nonzero.
\end{theorem}
\noindent\Realization{} \Kernel{} for the trace-derived dispersion, velocity
deficit inequality, exact massless set $\{0,\pi\}$, and uniform massive
discriminant bound; \Human{} for the implicit-differentiation identity in
\eqref{eq:groupvelocity} and for packaging these facts into the displayed
arccos branches.

\begin{proof}
Equation~\eqref{eq:bands} is the principal solution of
\eqref{eq:dispersion}.  A two-band degeneracy requires
$\sin\Omega=0$, equivalently
$|\cos q\cos\theta|=1$.  Each factor has absolute value at most one, so equality
forces $q,\theta\in\pi\mathbb Z$.  At $q=0$ one has
$\Omega_+=|\theta|$ on the stated branch; at $q=\pi$ one has
$\Omega_+=\pi-|\theta|$, whose circular distance to the negative band across
the quasienergy boundary is $2|\theta|$.  Implicit differentiation gives
\eqref{eq:groupvelocity}.  Finally,
\[
 \sin^2\Omega-\sin^2q\cos^2\theta
 =1-\cos^2\theta
 =\sin^2\theta\geq0,
\]
which proves the speed bound and its strict massive form.
\end{proof}

\begin{figure}[t]
\centering
\includegraphics[width=0.92\linewidth]{null_spinor_dirac_bands.pdf}
\caption{Exact quasienergy bands from
$\cos\Omega=\cos q\cos\theta$.  The left panel shows the massless cone at the
origin and its Floquet partner at quasienergy $\pi$.  A nonzero Pluecker area
opens both gaps (right).  The curves evaluate the exact formula; they are not a
fit or numerical simulation.  Generated by
\texttt{Scripts/sim/plot\_null\_spinor\_dirac\_bands.py}.}
\label{fig:bands}
\end{figure}

\paragraph{Doubling and pseudo-doubling verdict.}
The high-momentum $(\pi,\pi)$ cone is a Floquet zone-edge partner of the
low-energy cone.  It is not a second zero-quasienergy root, but it is a second
relativistic band-touching point in the full lattice spectrum.  Following the
terminology of Gupta and Short \cite{GuptaShortDoubling}, it is naturally
classified as a pseudo-doubler: a high-quasienergy mode with a Dirac-like local
dispersion.  The elementary walk therefore does \emph{not} establish a
no-doubling or no-pseudo-doubling theorem.  The bounded
continuum estimate in Section~\ref{sec:continuum} isolates fixed physical
momenta as $a\to0$, sending $q\simeq\pi$ to cutoff momentum; that decoupling is
not the same as removing the finite-lattice partner.

\paragraph{Chiral time frame.}
Let
\[
 U_{\mathrm s}(q,\theta)
 =e^{-\ii\theta\sigma_x/2}
  e^{-\ii q\sigma_z}
  e^{-\ii\theta\sigma_x/2}.
\]
It is unitarily equivalent to $U(q,\theta)$ and obeys the exact chiral relation
\begin{equation}\label{eq:chiral}
 \sigma_yU_{\mathrm s}(q,\theta)\sigma_y
 =U_{\mathrm s}(q,\theta)^\dagger.
\end{equation}
Thus the reciprocal pairing of quasienergies is structural, not an artifact of
the trace calculation.  Equation~\eqref{eq:chiral} is the symmetry statement
used here; a finer antiunitary or Floquet topological classification is left to
a separate analysis \cite{AsbothObuse}.

\paragraph{Robustness.}
If $z$ is perturbed to $z+\delta z$, then
\[
 \bigl||z+\delta z|-|z|\bigr|\leq|\delta z|.
\]
Therefore every perturbation with $|\delta z|<|z|$ remains off the collinear
gap-closing locus and retains a rest energy at least
$|z|-|\delta z|$.  The phase may rotate the internal basis, but by
\eqref{eq:phasecov} it cannot change a constant-coefficient spectrum.

\section{Uniform many-step Dirac control}
\label{sec:continuum}

The finite walk should approximate the continuous Dirac flow over a fixed
macroscopic time, not merely at one infinitesimal step.  For the real
representative of the phase-conjugacy class, write
\begin{equation}
 U(q,r)=e^{-\ii q\sigma_z}e^{-\ii r\sigma_x},
 \qquad
 H(k,\mu)=k\sigma_z+\mu\sigma_x,
 \qquad
 V_t(k,\mu)=e^{-\ii tH(k,\mu)}.
\end{equation}

\begin{theorem}[Uniform bounded-momentum product limit]\label{thm:rate}
For $K,M\geq0$, define
\begin{align}
 C_{\mathrm{box}}(K,M)
   &=2K^2+2M^2+KM^2+K^2M+KM,\\
 D_{\mathrm{box}}(K,M)
   &=4C_{\mathrm{box}}(K,M)
     +4(K+M)^2e^{K+M}.
\end{align}
If $|k|\leq K$, $|\mu|\leq M$, $n>0$, and $|t/n|\leq1$, then
\begin{equation}\label{eq:rate}
 \norm{
   U\!\left(k\frac{t}{n},\mu\frac{t}{n}\right)^n
   -V_t(k,\mu)}
 \leq D_{\mathrm{box}}(K,M)\frac{t^2}{n}.
\end{equation}
The bound is uniform on the displayed parameter box, and its right-hand side
tends to zero as $n\to\infty$.
\end{theorem}
\noindent\NewResult{} \Kernel{}

\begin{proof}[Proof architecture]
The one-step walk differs from $\id-\ii(t/n)H$ by an explicit quadratic
remainder.  The exact exponential differs from the same first-order matrix by
another quadratic remainder, giving
$\norm{U-V_{t/n}}\leq D(k,\mu)t^2/n^2$.  Both matrices are unitary, so the power
telescope has no norm-growth factor:
\[
 \norm{U^n-V_{t/n}^n}\leq n\norm{U-V_{t/n}}.
\]
The exact short-time flows compose to $V_t$, producing the $1/n$ rate.  Finally
$D(k,\mu)\leq D_{\mathrm{box}}(K,M)$ on the parameter box.  Lean checks the
entrywise remainder estimates, operator-norm conversion, unitarity, telescope,
flow composition, box inequality, and limit.
\end{proof}

\begin{figure}[t]
\centering
\includegraphics[width=0.88\linewidth]{null_spinor_continuum_rate.pdf}
\caption{Exact many-step operator errors for three nondegenerate parameter
choices.  All follow the proved $n^{-1}$ rate and remain below the explicit
compact-box envelope.  The points evaluate the exact finite matrices; the
envelope and convergence statement are kernel checked.}
\label{fig:continuumrate}
\end{figure}

This proof is a finite-dimensional first-order product-formula argument.  Its
constant is explicit rather than optimized; sharper general Trotter-error
theory provides context for that distinction \cite{ChildsTrotter}.

For $z=\mu u$ with $|u|=1$, the diagonal unitary $W_u$ commutes with $\Gamma$
and conjugates $B_\mu$ to $B_z$.  Unitary invariance of the operator norm
therefore transfers \eqref{eq:rate} to the oriented Pluecker walk.  This
one-line conjugacy corollary is \Human{} pending a combined declaration.

The theorem proves uniform convergence of the finite momentum symbol in
$1+1$ dimensions.  It is not yet an infinite-volume inverse-Fourier theorem,
a position-space PDE convergence theorem, or an interacting limit.  Those
quantifiers are separate problems, not implications hidden in the word
``continuum.''

\section{Finite position kernels and a controlled $3+1$ extension}
\label{sec:threeplusone}

The uniform estimates already imply a finite position-space statement.  Let
$\mathcal K$ be a finite momentum set, let $\chi_x(k)$ be phases with
$|\chi_x(k)|\leq1$, and define the synthesized matrix kernel
\[
 \mathcal S_x[A]:=\sum_{k\in\mathcal K}\chi_x(k)A(k).
\]
Applying the triangle inequality to Theorem~\ref{thm:rate} gives the
kernel-checked bound
\begin{equation}\label{eq:finitefourier}
 \norm{\mathcal S_x[U_n]-\mathcal S_x[V_t]}
 \leq |\mathcal K|D_{\mathrm{box}}(K,M)\frac{t^2}{n}.
\end{equation}
This controls every finite periodic Fourier kernel.  The cardinality factor is
the generic price of applying the triangle inequality pointwise.  It disappears
at the level of normalized $L^2$ wave packets.  For Mathlib's unnormalized
vector-valued transform $\mathcal F$ on $\mathbb Z/N\mathbb Z$, we prove
\begin{equation}\label{eq:plancherel}
 \sum_k
   \norm{(\mathcal F f)(k)}^2
   =N\sum_x
   \norm{f(x)}^2,
 \qquad
 \sum_x
   \norm{\mathcal F^{-1}g(x)}^2
   =\frac1N\sum_k
   \norm{g(k)}^2.
\end{equation}
Consequently, if a modewise error obeys
$\norm{\delta(k)}\leq\varepsilon\norm{c(k)}$, then
\begin{equation}\label{eq:l2wavepacket}
 \sum_x\norm{\mathcal F^{-1}\delta(x)}^2
 \leq \frac{\varepsilon^2}{N}
       \sum_k\norm{c(k)}^2.
\end{equation}
Both identities and the implication are \Classical{} \Kernel{}; a one-mode
packet is included as an exact normalization control.  We also prove the
operator-valued corollary: a uniform bound
$\norm{A(k)-B(k)}\leq\varepsilon$ implies
\eqref{eq:l2wavepacket} for
  $\delta(k)=(A(k)-B(k))c(k)$.  Applying that corollary to the local walk
  uses the exact plane-wave block stated below.  Iterating the transform in
  the three coordinates gives, for $N\geq1$, the kernel-checked product-torus identity
  \begin{equation}\label{eq:plancherel3}
    \sum_{x\in(\mathbb Z/N\mathbb Z)^3}
      \norm{\mathcal F_3^{-1}g(x)}^2
    =\frac{1}{N^3}
      \sum_{k\in(\mathbb Z/N\mathbb Z)^3}\norm{g(k)}^2.
  \end{equation}
  Hence, again for $N\geq1$, a uniform modewise operator error $\varepsilon$ gives the exact
  wave-packet coefficient $\varepsilon^2/N^3$, with no cardinality loss.  The
  momentum-sign bridge is also kernel-checked: for the positive-exponential
  convention the finite character block equals the analytic block at
  $k_j\varepsilon=-2\pi n_j/L$, including an exact quarter-zone $-\ii$ sign
  control.

  There is also a countable successor to the finite Plancherel statement.  If
  $e_n(k)$ is a countable family of mode errors such that
  $e_n(k)\to0$ for every $k$ and
  $\norm{e_n(k)}^2\leq b(k)$ for one summable envelope $b$, then
  \begin{equation}\label{eq:countablel2}
    \sum_k \norm{e_n(k)}^2\longrightarrow0.
  \end{equation}
  This is a direct vector-valued Tannery specialization, proved in Lean with a
  nonzero one-mode control.  It isolates the remaining physical obligation:
  construct the walk-specific square-summable envelope and identify the
  limiting multiplier with the Dirac PDE propagator under a continuum Fourier
  isometry.

  The corresponding measure-theoretic multiplier step is also closed.  If
  $f\in L^2$ and pointwise errors satisfy
  $\norm{e_n(k)}\leq\varepsilon_n\norm{f(k)}$ almost everywhere, where
  $0\leq\varepsilon_n$ and $\varepsilon_n\to0$, then
  \begin{equation}\label{eq:continuousl2multiplier}
    \norm{e_n}_{L^2}\longrightarrow0.
  \end{equation}
  The proof is a kernel-checked `eLpNorm` monotonicity-and-squeeze argument.
  The walk-specific compact-support application is now closed as well.  For any
  fixed $t\in\mathbb R$, suppose $f\in L^2$ and, almost everywhere on the support of $f$,
  $|k_x|,|k_y|,|k_z|\leq K$ and $|z|\leq M$, with $K,M\geq0$.  Applying the
  complex split symbol for $n+1$ steps gives an error multiplier $E_n$ satisfying
  \begin{equation}\label{eq:walkcompactl2}
    \norm{E_nf}_{L^2}\longrightarrow0,
    \qquad
    \norm{E_n(k)f(k)}
      \leq\frac{D_{K,M}t^2}{n+1}\norm{f(k)}.
  \end{equation}
  The displayed pointwise envelope holds for all sufficiently large $n$; the
  finite prefix is irrelevant to the sequence limit.
  This walk-specific composition is \NewResult{} \Kernel{}.  What remains is
  the lattice/continuum Fourier-isometry and identification of the limiting
  multiplier with the position-space Dirac PDE, not compact-support multiplier
  convergence.

The same architecture extends to four internal components.  The
successive-axis factorization itself is the standard construction used by
Mlodinow and Brun \cite{MlodinowBrun}; what is new here is the
Pluecker-derived complex coin, the exact finite-torus implementation with a
two-sided inverse and strict causal cone, and the kernel-checked all-zone
audit that follows.  Let Hermitian
matrices $\alpha_1,\alpha_2,\alpha_3,\beta\in M_4(\C)$ satisfy
\[
 \alpha_i^2=\beta^2=\id,
 \qquad \{\alpha_i,\alpha_j\}=0\ (i\neq j),
 \qquad \{\alpha_i,\beta\}=0,
\]
and set
\[
 H_4(\mathbf k,m)=\sum_{j=1}^3k_j\alpha_j+m\beta.
\]
The explicit matrices used in the artifact obey
\begin{equation}\label{eq:H4square}
 H_4(\mathbf k,m)^2=(|\mathbf k|^2+m^2)\id.
\end{equation}
For a substep $\varepsilon$, define
\begin{equation}\label{eq:3plus1step}
 U^{(4)}_\varepsilon(\mathbf k,m)
 =e^{-\ii\varepsilon k_x\alpha_1}
  e^{-\ii\varepsilon k_y\alpha_2}
  e^{-\ii\varepsilon k_z\alpha_3}
  e^{-\ii\varepsilon m\beta}.
\end{equation}
The order in \eqref{eq:3plus1step} is not cosmetic: the spatial generators do
not commute, and the artifact contains that noncommutation as a negative
control against a false second-order claim.

\begin{theorem}[Finite local $3+1$ walk and compact rate]
\label{thm:3plus1rate}
On a finite periodic three-dimensional position register, there is an explicit
successive-axis update built from one-site conditional shifts and the mass coin
in \eqref{eq:3plus1step}.  It preserves the finite inner product exactly, has
an explicit two-sided inverse, and reads only sites reached by one conditional
source shift along each axis in one cycle.  Its product plane-wave sectors are
invariant, and on each sector the update is
exactly the ordered axis-$0$/axis-$1$/axis-$2$/mass matrix block built from
finite roots of unity.  For the positive-exponential plane-wave convention,
that block is \eqref{eq:3plus1step} when
\begin{equation}\label{eq:latticemomentum}
 k_j\varepsilon=-\frac{2\pi n_j}{L}.
\end{equation}
If
$|k_x|,|k_y|,|k_z|\leq K$, $|m|\leq M$, $n>0$, and $|t/n|\leq1$, then
\begin{equation}\label{eq:3plus1ratebound}
 \norm{\left(U^{(4)}_{t/n}(\mathbf k,m)\right)^n
       -e^{-\ii tH_4(\mathbf k,m)}}
 \leq 16(3K+M)^2e^{3K+M}\frac{t^2}{n}.
\end{equation}
Every finite Fourier synthesis obeys the same bound multiplied by its number
of modes, while the normalized product-torus $L^2$ wave-packet error obeys
the squared bound divided by $N^3$.  Moreover, let $B_4(z)$ be the
phase-retaining rest operator and let $U^{(4)}_\varepsilon(\mathbf k,z)$ be the
chiral conjugate of the displayed step at $m=|z|$.  Its exact comparison flow
is generated by
\[
 H_4(\mathbf k,z)=\sum_{j=1}^3k_j\alpha_j+B_4(z),
 \qquad
 H_4(\mathbf k,z)^2=(|\mathbf k|^2+|z|^2)\id,
\]
and whenever $|z|\leq M$ it satisfies the same bound
\begin{equation}\label{eq:complex3plus1rate}
 \norm{\left(U^{(4)}_{t/n}(\mathbf k,z)\right)^n
       -e^{-\ii tH_4(\mathbf k,z)}}
 \leq 16(3K+M)^2e^{3K+M}\frac{t^2}{n}.
\end{equation}
\end{theorem}
\noindent\Realization{} \Kernel{} for exact local norm preservation, Clifford
square, finite-character plane-wave block, compact rate, finite synthesis
bound, and the trigonometric/sign conversion
\eqref{eq:latticemomentum} to the analytic symbol; \NewResult{} \Kernel{} for
the complex-phase transfer \eqref{eq:complex3plus1rate} and the iterated
three-torus normalization \eqref{eq:plancherel3}.

\begin{proof}[Proof architecture]
Each $\alpha_j$ has eigenvalues $\pm1$.  An explicit unitary eigenbasis turns
the corresponding diagonal one-site conditional shift into the physical
Clifford factor; composing those shifts with the mass coin proves exact norm
preservation.  Direct substitution into every product plane wave extracts one
unit-modulus character per conditional shift and proves the exact ordered
finite block without an asymptotic argument.  The closed form
$e^{-\ii q\alpha_j}=\cos q\,\id-\ii\sin q\,\alpha_j$ identifies the ordered
analytic matrix symbol after the sign convention in
\eqref{eq:latticemomentum}.  A four-factor quadratic remainder estimate, exact
unitarity, and the unitary power telescope then give
\eqref{eq:3plus1ratebound}.  Applying the finite synthesis theorem gives its
position-kernel corollary.
\end{proof}

For $z\neq0$, the phase transfer is exact.  If $W_z$ is the chiral unitary at
the polar angle of $z$, then
\[
 W_zH_4(\mathbf k,|z|)W_z^\dagger=H_4(\mathbf k,z),
 \qquad
 W_ze^{-\ii tH_4(\mathbf k,|z|)}W_z^\dagger
   =e^{-\ii tH_4(\mathbf k,z)}.
\]
The operator norm cannot increase under this conjugation, so the real-mass
estimate transfers without changing its constant.

\subsection*{The $3+1$ high-symmetry verdict}

The continuum tangent is exactly isotropic:
$H_4(\mathbf k,0)^2=|\mathbf k|^2\id$.  The finite cubic walk is not globally
doubler-free.  At the eight corners
$(q_x,q_y,q_z)\in\{0,\pi\}^3$, its massless symbol is
\begin{equation}\label{eq:cornerparity}
 U^{(4)}(q_x,q_y,q_z,0)=(-1)^{r}\id,
 \qquad
 r=\#\{j:q_j=\pi\}.
\end{equation}
Thus $(\pi,\pi,0)$, $(\pi,0,\pi)$, and $(0,\pi,\pi)$ are distinct momenta
 with exactly the same zero-quasienergy symbol as the origin.  Equation
 \eqref{eq:cornerparity}, all three distinctness witnesses, and radial isotropy
 of the continuum square are \Realization{} \Kernel{}.  This is an actual
 finite-regulator obstruction, not a caveat inferred from a plot.

The massive walk has a stronger interior obstruction.  At body-center momentum
$\mathbf q_B=(\pi/2,\pi/2,\pi/2)$, let
$u_\theta=\cos\theta+\ii\sin\theta$ and
\[
 v_+(\theta)=(-u_\theta,0,1,0)^T,
 \qquad
 v_-(\theta)=(u_\theta,0,1,0)^T.
\]
For every real mass angle $\theta$, the exact finite step obeys
\begin{equation}\label{eq:bodycentermodes}
 U^{(4)}(\mathbf q_B,\theta)v_+(\theta)=v_+(\theta),
 \qquad
 U^{(4)}(\mathbf q_B,\theta)v_-(\theta)=-v_-(\theta),
 \qquad v_\pm(\theta)\neq0.
\end{equation}
Equivalently, the body-center step is a non-scalar involution for all
$\theta$.  It therefore retains exact quasienergies $0$ and $\pi$ even after
the mass coin is turned on: the present cubic regulator admits no uniform
Floquet gap (a body-center crossing at $0$ and $\pi$ persists for every
mass angle).
The closed body-center matrix, involution identity, and both explicit
eigenmodes are \NoGo{} \Kernel{}.

The obstruction is also visible in an exact all-zone criterion.  Write
$x=\cos q_x$, $y=\cos q_y$, $w=\cos q_z$, and $c=\cos\theta$, and define
\begin{align}
 S={}&4c^2x^2y^2w^2-2c^2x^2y^2-2c^2x^2w^2+c^2x^2
      -2c^2y^2w^2+c^2y^2+c^2w^2 \notag\\
   &{}-2x^2y^2w^2+x^2y^2+x^2w^2+y^2w^2,\\
 P_0&=S-2cxyw,
 \qquad
 P_\pi=S+2cxyw.
\end{align}
For the live ordered successive-axis symbol, the finite determinant expansion
is now exact at every real momentum and mass angle:
\begin{equation}\label{eq:fullblochdet}
 \det\!\left(U^{(4)}-\id\right)=4P_0,
 \qquad
 \det\!\left(U^{(4)}+\id\right)=4P_\pi.
\end{equation}
Thus $P_0=0$ and $P_\pi=0$ are respectively necessary and sufficient for
zero- and $\pi$-quasienergy modes.  At body center $x=y=w=0$, both polynomials
vanish for every $\theta$, recovering \eqref{eq:bodycentermodes}.  The two
determinant identities, their explicit $4\times4$ expansion, and the
definitional bridge to the walk used in the continuum estimate are
\NoGo{} \Kernel{}.

The zero sets are now classified on the principal massive branch.  If
$0<|c|<1$ and $|x|,|y|,|w|\leq1$, then
\begin{equation}\label{eq:fullblochzeroset}
 P_0=0 \quad\Longleftrightarrow\quad x=y=w=0
 \quad\Longleftrightarrow\quad P_\pi=0.
\end{equation}
The proof uses the exact decomposition
\begin{equation}
 P_0=(cx-yw)^2+
 \bigl(y^2+w^2-2y^2w^2\bigr)
 \bigl(c^2+x^2(1-2c^2)\bigr),
\end{equation}
whose second factor is strictly positive and whose first prefactor is
nonnegative on this branch.  Equality then forces all three momentum cosines
to vanish; the $P_\pi$ statement follows by $c\mapsto-c$.  Thus the simultaneous
body-center cosine locus is the complete zero- and $\pi$-quasienergy crossing
set for every principal massive angle, not merely a sampled obstruction.
The hypotheses are load-bearing: at $c=0$, additional zeros occur, including
$(x,y,w)=(1,0,0)$.  The classification and this boundary control are
\NoGo{} \Kernel{}.  They classify the failure of the present regulator; they
do not make it alias-free.

\paragraph{A degree-one stationary-amplitude no-go.}
The most economical proposed repair is to add a stay-put channel to each
nearest-neighbor factor.  The exact ansatz is
\begin{equation}\label{eq:stationarylaurent}
 F(q)=e^{\ii q}A+B+e^{-\ii q}C.
\end{equation}
Let $M$ be any Hermitian involution.  If $F(q)$ is unitary for every real $q$,
$F(0)=\id$, and $F'(0)=-\ii M$, then
\begin{equation}\label{eq:stationarynogo}
 B=0.
\end{equation}
The proof evaluates the exact unitarity residual at $q=\pm\pi/2$; adding the
two resulting matrix identities forces $B^\dagger B=0$.  Specializing
$M=\alpha_j$ rules out a nonzero stationary amplitude in a degree-one factor
with the full Dirac tangent, even before zone-edge separation is imposed.
All three specializations to the live $3+1$ generators are separately
kernel-checked.  This is \NoGo{} \Kernel{}.  A nonvacuity control shows that a partial-channel
Hermitian tangent with a stationary kernel does admit an exactly unitary
degree-one witness with $B\neq0$ and a non-scalar zone edge.  That witness does
not retain the full Dirac tangent.  Thus stationary amplitudes remain viable
only after a genuine change of architecture, such as a larger internal cell,
longer Laurent range, multistep factorization, or a non-involutory tangent.

The last route is not hypothetical.  Gupta and Short
\cite{GuptaShortDoubling} construct doubler- and pseudo-doubler-free
nearest-neighbor one-step unitary walks, in $1+1$ and $3+1$ dimensions, whose
enabling resource is exactly a nonzero stay-put amplitude
$\gamma_0=\id-\gamma_+-\gamma_-$ built from tilted, non-orthogonal velocity
projectors (their Eqs.~(29)--(30) and~(37)).  Their step is exactly of the
degree-one Laurent form \eqref{eq:stationarylaurent}, unitary at every
momentum and normalized at the origin, with stationary block $\gamma_0$.
The no-go \eqref{eq:stationarynogo} therefore yields a corollary about
their family: its tangent generator, though automatically Hermitian for a
unitary origin-normalized step, cannot be an involution ($M^2
\neq\id$).
The stay-put freedom that removes the doublers is purchased
exactly by relinquishing the involutory unit-speed Dirac tangent, and our
partial-channel escape witness realizes the same mechanism in kernel-checked
form.  Their $3+1$ family moreover retains extraneous Weyl-like low-energy
solutions at isolated momenta $\pm\mathbf q(\theta)$, which the authors
suggest a more general higher-dimensional construction might remove (their
Appendix~F), so a strictly local, exactly unitary, fully alias-free $3+1$
walk with the complete involutory Dirac tangent remains open in both
programs.

The obstruction is stronger for a separable three-axis product.  Once
$B_x=B_y=B_z=0$, origin normalization forces each factor to obey
$F_j(\pi)=-\id$.  Therefore, for any momentum-independent onsite coin $Q$,
including the complex Pluecker mass coin,
\begin{equation}\label{eq:genericcorneralias}
 U_Q(0,0,0)=U_Q(\pi,\pi,0)=U_Q(\pi,0,\pi)=U_Q(0,\pi,\pi).
\end{equation}
This all-coins alias theorem is \NoGo{} \Kernel{}.  It proves that changing
the onsite mass data cannot repair the cubic partners inside the
four-channel, range-one, single-factor-per-axis architecture.  Extra range,
coupled substeps, or enlarged/ancillary cells are not optional embellishments;
at least one is structurally required.

An attempted commutator--Wilson repair sharpens that requirement.  Consider
the exactly unitary one-axis gate
\begin{equation}
 W_r(q)=F_{\beta}(rq)F_{\gamma_5}(rq)
        F_{\beta}(-rq)F_{\gamma_5}(-rq),
 \qquad
 F_A(s)=\cos s\,\id-\ii\sin s\,A.
\end{equation}
Its linear term cancels, so at the Bloch-matrix level it appears to supply a
second-order Wilson-like correction without disturbing the Dirac tangent.
However, in the fixed one-cell ansatz strict finite Laurent range requires the
frequency $r$ to be an integer.  For every integer $r$ the exact zone-edge
identity is
\begin{equation}
 W_r(\pi)=\id,
\end{equation}
and all three even-corner aliases in \eqref{eq:genericcorneralias} remain
exactly equal to the origin.  Conversely, the noninteger trigonometric values
that lift selected corners are not $2\pi$-periodic finite Laurent polynomials
in $e^{\ii q}$ and therefore do not define the claimed strict local one-cell
QCA.  This family-scoped locality/de-aliasing incompatibility, including the
range-one control, is \NoGo{} \Kernel{}.  It does not exclude a larger cell or
a genuinely different finite product; it prevents a nonlocal Bloch control
from being mislabeled as that successor.

The matching cross-architecture algebraic invariant is now exact in one
Laurent variable.  The determinant of every invertible finite Laurent matrix
is a nonzero scalar times a uniquely exponented monomial, and that exponent is
additive under composition, vanishes on the identity, and equals $n$ on the
one-channel shift $T^n$.  This is the algebraic precursor of the
translation-invariant one-particle quantum-walk index in GNVW Proposition~5
\cite{GNVW}; identifying the two still requires a complex-unitary
Laurent/Fourier representation theorem and an explicit shift-sign convention.
It is not the positive-rational many-body GNVW QCA index, an entropy-flow
theorem, or a no-doubling result.

There is already a constructive Hamiltonian-level route around this no-go.
Using the same live Clifford generators, define
\begin{align}
 M_W(\mathbf q)&=m+r\sum_{j=1}^{3}(1-\cos q_j),\\
 H_W(\mathbf q)&=\sum_{j=1}^{3}\sin q_j\,\alpha_j+M_W(\mathbf q)\beta.
\end{align}
The kernel-checked Wilson identity is
\begin{equation}\label{eq:wilsonsquare}
 H_W(\mathbf q)^2=
 \left(\sum_{j=1}^{3}\sin^2q_j+M_W(\mathbf q)^2\right)\id.
\end{equation}
For $m>0$ and $r\geq0$, the scalar coefficient is bounded below by $m^2$.
For $m=0$ and $r>0$, it vanishes if and only if
$\cos q_1=\cos q_2=\cos q_3=1$; hence all non-origin cubic corners are lifted.
For example, the $(\pi,0,0)$ energy squared is exactly $4r^2$.  This closes
the spectral design problem at the nearest-neighbor Hamiltonian level while
preserving the Dirac tangent.  It does \emph{not} yet close the stricter QCA
problem: $\exp(-\ii tH_W)$ is exactly unitary, but its finite-time exponential
is not asserted to be a finite Laurent polynomial or a strictly finite-range
one-step update.

\paragraph{Exact temporal blocking and nonclosure.}
Define the two-step blocking map by $\mathcal B(U)=U^2$.  The mass-only
Pluecker subgroup closes exactly,
\[
 \mathcal B\bigl(C_z(\varepsilon)\bigr)=C_z(2\varepsilon),
\]
but the ordered spatial split family does not close by doubling its scalar
step.  At the exact $x$--$y$ quarter-turn witness,
\begin{equation}\label{eq:rgcounterexample}
 \mathcal B\bigl(U_{xy}(\pi/2)\bigr)=-\id,
 \qquad U_{xy}(\pi)=\id.
\end{equation}
The smallest multiplicatively closed enlargement is the submonoid generated
by all bare split steps.  These statements are \NewResult{} \Kernel{}: temporal
blocking generates ordered word/coupling data rather than merely a
renormalized copy of the original scalar parameters.  Spatial decimation is a
separate, stronger RG problem.

\paragraph{A proved boundary on scale selection.}
For the finite positive Pluecker action developed here, simultaneous scaling
of both primitive spinors by $t$ gives the exact law
\begin{equation}\label{eq:quarticscalenogo}
 S(t\psi,t\phi;x)=t^4S(\psi,\phi;x).
\end{equation}
Whenever the unscaled action is positive, its radial equation
$4t^3S(\psi,\phi;x)=0$ is equivalent to $t=0$.  The orthogonal-spinor,
positive-quartet control has $S=1/2$, so the obstruction is nonvacuous.
The formal result is now more general.  For any positive finite action with
radial profile $S(tx)=t^pS(x)$ and natural degree $p>0$, the derivative at unit
scale is $pS(x)>0$, so the configuration cannot be stationary.  At $p=0$ the
radial direction is exactly flat and cannot select an isolated scale.  This is
\NoGo{} \Kernel{}: no positive fixed finite homogeneous action can
spontaneously select a nonzero scale.  A running/refining family, competing
homogeneity, constraint, ensemble scale, or other dimensionful input must be
displayed.
The exact positive control $V_c(t)=(t^2-c)^2$ has stationary points precisely
at $t=0$ or $t^2=c$; for supplied $c=1$, the nonzero points $t=\pm1$ are global
minima.  Selection is therefore possible, but its magnitude is exactly the
new input.  This control is also \Realization{} \Kernel{}.

\paragraph{Pre-registered high-momentum benchmark.}
Before evaluation we fixed $z=3+4\ii$, $\varepsilon=0.2$, Wilson $r=1$, and a
$10^{-10}$ kill threshold.  The oracle then found exactly three non-origin
even-corner aliases (maximum residual $2.30\times10^{-16}$), both body-center
$\pm1$ modes (residuals below $2.24\times10^{-16}$), no non-origin Wilson zero,
and minimum Wilson corner gap $2$.  With $r=0$, all eight corners became
gapless as required by the negative control.  No parameter was fitted.  This
numerical result validates the implementation of the theorem chain; the Lean
identities, not the benchmark, carry the proof.

\begin{figure}[t]
\centering
\includegraphics[width=0.98\linewidth]{null_spinor_3plus1_brillouin.pdf}
\caption{High-symmetry audit of the ordered $3+1$ walk.  The small-momentum
tangent is radially isotropic, but the exact massless lattice symbol has four
zero-quasienergy corners: the origin and three even-parity aliases.  The marked
body-center crossing in the massive panel persists for every mass angle by
Eq.~\eqref{eq:bodycentermodes}.  Curves evaluate the exact finite matrices; the
corner classification and body-center eigenmodes are kernel checked.}
\label{fig:3plus1bz}
\end{figure}

\paragraph{Comparison with the tetrahedral walk.}
Nzongani et al.~\cite{NzonganiTetra} place four amplitudes on tetrahedral
facets, compose facet/strict/weakly local shifts with directional coins, and
insert the assigned coin $e^{-\ii\varepsilon m\gamma^0}$; their expansion
recovers the $3+1$ Dirac equation to first order.  Our register and finite
update are different: the present construction is a cubic successive-axis
product with a Pluecker-derived complex mass coin and an explicit many-step
rate.  A constant onsite basis change cannot decide away this difference:
unitary conjugacy preserves equality of symbols at distinct momenta, so every
 constant onsite conjugate of our walk inherits the three aliases above and
 the body-center $\{+1,-1\}$ spectrum.  Any equivalence with a tetrahedral update
 must therefore reproduce those obstructions
or use a non-onsite blocking, changed unit cell, or coarse-graining map.  A
full spectral audit of the tetrahedral symbol is needed before choosing among
those possibilities.

\paragraph{Small-step-sensitive $3+1$ rate.}
The live four-factor symbol now has a sharper fixed-time theorem.  With
$B_4=|k_x|+|k_y|+|k_z|+|m|$ and $n>0$,
\[
 \left\|U(t/n)^n-e^{-\ii tH}\right\|
 \leq {2B_4^2t^2\over n}\exp\!\left({|t|B_4\over n}\right).
\]
The proof keeps the small-step exponent through both the split-product
remainder and the unitary telescope; unlike the earlier compact-box envelope,
it does not replace this factor by $\exp(B_4)$.  The resulting growing-window
theorem takes $n=K^4$ and permits arbitrary momentum and mass sequences with
$B_4\leq K$; their modewise split-versus-exact error is bounded by
$2t^2K^{-2}\exp(|t|K^{-3})$ and tends to zero.  A sharp pure-momentum boundary
fixture saturates $B_4=K$, so this is not a fixed-mode restatement.  Shannon
sampling/interpolation into physical space and the position-space $L^2$ Dirac
PDE limit remain Paper~D.

\section{Formal verification and negative controls}

The formal development uses Lean~4.28.0 and Mathlib \cite{Lean4,Mathlib}.  Its
spinor and physics conventions were also cross-checked against the public
PhysLean corpus \cite{PhysLean}.  Formalized mathematical physics at
substantial scale is becoming practical: the free bosonic quantum field
theory in four-dimensional Euclidean spacetime, satisfying the Glimm--Jaffe
axioms, was recently constructed end to end in Lean~4
\cite{DouglasQFTLean}.  The present artifact is complementary in scope:
exact finite discrete dynamics, with the physical input --- the rest
operator itself --- derived inside the formalization rather than
axiomatized.  The trusted
Cauchy--Binet identity lives in the project core; the newer operator and walk
results remain in the draft namespace while their publication interfaces are
reviewed.  Each headline declaration below has a build-enforced assumption
footprint contained in
$\{\mathtt{propext},\mathtt{Classical.choice},\mathtt{Quot.sound}\}$.

\begin{center}
\small
\begin{tabularx}{\linewidth}{@{}p{0.29\linewidth}p{0.15\linewidth}X@{}}
\toprule
Declaration & Status & Mathematical payload\\
\midrule
\texttt{two\_edge\_plucker\_mass\_identity}
 & \Classical{} \Kernel{} & $\det P=|\psi\wedge\phi|^2$\\
\texttt{fin\_bundle\_plucker\_mass\_identity}
 & \Classical{} \Kernel{} & finite-family Cauchy--Binet area sum\\
\texttt{hermitian\_determinant\_uniqueness}
 & \Classical{} \Kernel{} & determinant is the unique quadratic null-boundary form up to scale\\
\texttt{pair\_diracSymbol\_sq}
 & \NewResult{} \Kernel{} & $H_z(k)^2=(k^2+\det P)\id$ for arbitrary pairs\\
\texttt{massOperator\_eq\_zero\_iff\_collinear}
 & \NewResult{} \Kernel{} & exact gap-closing locus\\
\texttt{phase\_covariance}
 & \NewResult{} \Kernel{} & oriented phase acts by conjugacy\\
\texttt{massOperator\_sl2\_invariant}
 & \NewResult{} \Kernel{} & invariance under the defining $\mathrm{SL}(2,\C)$ action\\
\texttt{oddHermitian\_eq\_massOperator}
 & \NewResult{} \Kernel{} & complete normal form for odd Hermitian rest operators\\
\texttt{massCoin\_unitary\_group}
 & \NewResult{} \Kernel{} & exact unitary rest evolution and group law\\
\texttt{unitary\_kernel\_eq\_scaled\_checkerboard}
 & \Realization{} \Kernel{} & exact $\cos^n$ scaling of the two-weight recursion; operator/history identification is supplied by separate transfer-matrix theorems\\
\texttt{exact\_quantum\_walk\_dispersion\_verdict}
 & \Realization{} \Kernel{} & concrete trace identity, trace-to-quasienergy bridge, unitarity, determinant one, and rest eigenphases\\
\texttt{massive\_implies\_subluminal}
 & \Realization{} \Kernel{} & strict full-zone velocity-deficit inequality for massive angle; derivative identity is \Human{}\\
\texttt{massless\_gapless\_classification}
 & \Realization{} \Kernel{} & exact $\{0,\pi\}$ massless band-touching set\\
\texttt{pathAmplitude\_eq\_corner\_power}
 & \Realization{} \Kernel{} & exact real-mass Gaussian-rational corner-power weight; complex orientation is carried by phase conjugacy\\
\texttt{complexPlueckerPathSum\_eq\_transfer\_pow}
 & \NewResult{} \Kernel{} & direct directed-edge history sum using the actual complex coin entries $z$ and $\bar z$\\
\texttt{fixed\_time\_many\_step\_bound}
 & \NewResult{} \Kernel{} & fixed-parameter $D(k,\mu)t^2/n$ estimate\\
\texttt{fixed\_time\_many\_step\_bound\_on\_box}
 & \NewResult{} \Kernel{} & uniform bounded-parameter estimate\\
\texttt{localStep\_preserves\_norm}
 & \Realization{} \Kernel{} & exact finite-torus $3+1$ norm preservation\\
\texttt{inverseSuccessiveWalk\_after}, \texttt{successiveWalk\_after\_inverse}
 & \Realization{} \Kernel{} & explicit two-sided inverse of the complete local cycle\\
\texttt{successiveWalk\_zero\_of\_causalSource\_zero}
 & \Realization{} \Kernel{} & strict one-shift-per-axis causal cone\\
\texttt{Finite3Plus1FourierBridge.localStep\_mode}
 & \Realization{} \Kernel{} & exact finite-character block on every product plane wave\\
\texttt{finiteLocalSymbol\_eq\_analytic\_neg}
 & \Realization{} \Kernel{} & exact finite-to-analytic momentum-sign conversion\\
\texttt{quarter\_zone\_sign\_control}
 & \Realization{} \Kernel{} & nondegenerate $-\ii$ convention fixture\\
\texttt{Pluecker3Plus1ComplexMass.H4\_sq}
 & \NewResult{} \Kernel{} & phase-retaining $3+1$ mass gives $H_4^2=(|\mathbf k|^2+|z|^2)\id$\\
\texttt{complex\_phase\_covariance}
 & \NewResult{} \Kernel{} & chiral covariance of the complex four-component rest operator\\
\texttt{massCoin4\_unitary\_group}
 & \NewResult{} \Kernel{} & derived complex mass coin is unitary with exact group law\\
\texttt{variableComplexLocalStep\_preserves\_norm}
 & \NewResult{} \Kernel{} & arbitrary position-dependent Pluecker profiles, including zero sites, evolve unitarily\\
\texttt{variableComplexLocalStep\_zero\_of\_causalSource\_zero}
 & \NewResult{} \Kernel{} & strict causal cone for the inhomogeneous complex-Pluecker walk\\
\texttt{local\_phase\_connection\_covariance}
 & \NewResult{} \Kernel{} & site-dependent phase induces exact endpoint links and rotates $z(x)$ pointwise\\
\texttt{gaugedCliffordAxisShift\_constant}
 & \NewResult{} \Kernel{} & constant phase produces no endpoint link\\
\texttt{global\_real\_lift\_forces\_zero\_winding}
 & \NoGo{} \Kernel{} & a global real phase lift has zero winding; nonzero index needs branch/transition data\\
\texttt{threeLinkUnitWinding\_not\_global\_lift}
 & \NewResult{} \Kernel{} & an explicit periodic link field has winding one and cannot be integrated to a global real vertex phase\\
\texttt{totalTurning\_int\_multiple}
 & \NewResult{} \Kernel{} & the link increments derived from any nowhere-zero Pluecker field (principal arguments of neighbor ratios) sum to an exact $2\pi$-integer winding\\
\texttt{windingOneField\_totalTurning},
\texttt{windingOneSpinors\_pluecker}
 & \NewResult{} \Kernel{} & an explicit four-site field has derived total turning exactly $2\pi$ and is the Pluecker coordinate of a primitive spinor pair field\\
\texttt{totalTurning\_eq\_zero\_of\_global\_lift},
\texttt{totalTurning\_const},
\texttt{linkIncrement\_global\_phase}
 & \NewResult{} \Kernel{} & derived increments matching a global real lift force zero turning; constant-field control; global chiral phase changes no derived increment\\
\texttt{ComplexPlueckerLocalWalk.complexLocalStep\_mode}
 & \NewResult{} \Kernel{} & phase-retaining mass coin inserted into the exact local Fourier block\\
\texttt{complex\_fixed\_time\_many\_step\_bound\_on\_box}
 & \NewResult{} \Kernel{} & complex Pluecker $3+1$ walk inherits the uniform $O(1/n)$ rate\\
\texttt{complexSplitStep\_eq\_explicit}
 & \NewResult{} \Kernel{} & for $z\neq0$, the conjugacy-defined rate step equals the implemented complex mass-coin product\\
\texttt{localSymbol\_many\_step\_bound\_on\_box}
 & \NewResult{} \Kernel{} & ordered local $3+1$ symbol has uniform $O(1/n)$ rate\\
\texttt{threePlusOne\_finite\_kernel\_bound}
 & \NewResult{} \Kernel{} & finite $3+1$ Fourier-kernel convergence\\
\texttt{dft\_energy}, \texttt{invDFT\_energy}
 & \Classical{} \Kernel{} & exact vector-valued finite Plancherel normalization\\
\texttt{inverseDFT\_wavepacket\_error}
 & \Classical{} \Kernel{} & modewise relative error implies normalized finite $L^2$ control\\
\texttt{inverseDFT\_operator\_wavepacket\_error}
 & \Classical{} \Kernel{} & uniform operator error controls every finite wave packet\\
\texttt{invDFT3\_energy}
 & \Classical{} \Kernel{} & exact inverse-transform energy factor $N^{-3}$ on the three-torus\\
\texttt{inverseDFT3\_operator\_wavepacket\_error}
 & \Classical{} \Kernel{} & exact normalized three-axis operator wave-packet bound\\
\texttt{complex\_walk\_inverseDFT3\_wavepacket\_error\_on\_box}
 & \NewResult{} \Kernel{} & walk-specific complex $3+1$ finite-torus $L^2$ rate\\
\texttt{countable\_l2\_error\_tendsto}
 & \Classical{} \Kernel{} & dominated mode convergence implies countable spectral $L^2$ convergence\\
\texttt{eLpNorm\_two\_tendsto\_zero\_of\_uniform\_relative\_bound}
 & \Classical{} \Kernel{} & uniform relative multiplier error implies measure-theoretic $L^2$ convergence\\
\texttt{walk\_error\_eLpNorm\_tendsto\_zero\_on\_compact\_support\_any\_time}
 & \NewResult{} \Kernel{} & at every fixed time, the complex $3+1$ walk converges in momentum-space $L^2$ on bounded support\\
\texttt{tail\_bulk\_split}
 & \Classical{} \Kernel{} & for any measurable band, restricted approximation error is bounded by its in-band error plus the complementary tail\\
\texttt{uv\_tail\_tendsto\_zero}
 & \Classical{} \Kernel{} & for arbitrary $L^2$ data, mass outside measurable monotone bands vanishes as the bands exhaust frequency space\\
\texttt{band\_approx\_tendsto\_zero}
 & \Classical{} \Kernel{} & abstract vanishing in-band error plus measurable band exhaustion implies global $L^2$ convergence; concrete sampling remains open\\
\texttt{massless\_corner\_parity\_classification}
 & \Realization{} \Kernel{} & exact $\{0,\pi\}^3$ matrix parity classification\\
\texttt{massless\_corner\_quasienergy\_classification}
 & \Realization{} \Kernel{} & exact corner Floquet-quasienergy classification\\
\texttt{origin\_alias\_pi\_pi\_zero}
 & \NoGo{} \Kernel{} & explicit distinct zero-quasienergy corner alias\\
\texttt{body\_center\_walk\_sq}
 & \NoGo{} \Kernel{} & massive body-center step is an involution for every mass angle\\
\texttt{body\_center\_plus\_mode\_eigen},
\texttt{body\_center\_minus\_mode\_eigen}
 & \NoGo{} \Kernel{} & explicit exact $+1$ and $-1$ body-center eigenmodes\\
\texttt{det\_splitStep\_sub\_one},
\texttt{det\_splitStep\_add\_one}
 & \NoGo{} \Kernel{} & exact all-momentum determinant polynomials for zero- and $\pi$-quasienergy modes of the live ordered walk\\
\texttt{zeroModePolynomial\_eq\_zero\_iff},
\texttt{piModePolynomial\_eq\_zero\_iff}
 & \NoGo{} \Kernel{} & complete principal-massive zero/$\pi$ crossing classification: both vanish exactly when all three momentum cosines vanish\\
\texttt{quarter\_turn\_mass\_extra\_zero}
 & \NoGo{} \Kernel{} & exact $\cos\theta=0$ boundary control with an additional non-body-center cosine zero\\
\texttt{stationary\_forces\_zero}, \texttt{dirac\_no\_go}
 & \NoGo{} \Kernel{} & origin normalization, exact unitarity, and a full involutory tangent exclude a degree-one stationary amplitude\\
\texttt{alpha1/2/3\_stationary\_forces\_zero}
 & \NoGo{} \Kernel{} & the no-go applies to each live $3+1$ axis generator\\
\texttt{live\_degree\_one\_factorized\_lower\_bound}
 & \NoGo{} \Kernel{} & every even-parity corner aliases the origin after any momentum-independent onsite coin\\
\texttt{wilsonAxis\_pi\_int},
\texttt{successor\_even\_corners\_int}
 & \NoGo{} \Kernel{} & every strict integer-frequency member of the commutator--Wilson ansatz has trivial zone-edge gate and retains all three aliases\\
\texttt{no\_integer\_pythagorean\_zone\_edge},
\texttt{range\_one\_alias\_control}
 & \NoGo{} \Kernel{} & noninteger de-aliasing certificate cannot occur at finite integer harmonic range; explicit range-one nonvacuity control\\
\texttt{relaxed\_witness}
 & \Realization{} \Kernel{} & a partial-channel tangent with stationary kernel admits a nonzero stationary amplitude\\
\texttt{WilsonDiracRegulator.H\_sq}
 & \NewResult{} \Kernel{} & exact Wilson--Dirac scalar square in the live Clifford basis\\
\texttt{massless\_energy\_eq\_zero\_iff}
 & \NewResult{} \Kernel{} & positive Wilson parameter removes every non-origin cubic corner at the Hamiltonian level\\
\texttt{bare\_mass\_sq\_le\_energySq}
 & \NewResult{} \Kernel{} & a positive bare mass supplies a uniform global Hamiltonian gap\\
\texttt{no\_constant\_onsite\_equivalence\_of\_corner\_separation}
 & \NoGo{} \Kernel{} & onsite conjugacy cannot erase a corner alias\\
\texttt{massCoin4\_block2\_closes},
\texttt{xy\_blocking\_nonclosure\_control}
 & \NewResult{} \Kernel{} & mass-only temporal blocking closes, while the ordered split family fails exact scalar-parameter closure\\
\texttt{no\_stationary\_positive\_homogeneous\_action},
\texttt{no\_nonzero\_stationary\_scale}
 & \NoGo{} \Kernel{} & no positive fixed finite homogeneous action of positive natural degree selects a stationary nonzero radial scale; the Pluecker action is the quartic instance\\
\texttt{create\_annihilate\_same\_add},
\texttt{create\_create\_anticomm}
 & \NewResult{} \Kernel{} & generic finite occupation-basis creation and annihilation maps satisfy the CAR\\
\texttt{Gamma\_vac}, \texttt{Gamma\_apply\_singleton},
\texttt{Gamma\_number}, \texttt{Gamma\_parity}
 & \NewResult{} \Kernel{} & determinant-minor second quantization fixes vacuum, equals $U$ on one-particle states, and preserves number and parity\\
\texttt{gamma\_create\_covariance}
 & \NewResult{} \Kernel{} & the determinant-minor lift sends creation in an input mode to the exact $U$-weighted sum of output-mode creations\\
\texttt{Gamma\_one}, \texttt{gammaEntry\_mul},
\texttt{Gamma\_mul}
 & \NewResult{} \Kernel{} & exact identity and composition laws for the determinant-minor exterior lift, including ordered-minor Cauchy--Binet\\
\texttt{plueckerQuartic\_forward\_amplitude},
\texttt{witness\_pair\_orientation\_sensitive}
 & \NewResult{} \Kernel{} & a quartic two-particle transfer reads the primitive wedge $z$ and distinguishes $3+4\ii$ from its conjugate\\
\texttt{pairKick\_add}, \texttt{pairKick\_smul},
\texttt{witnessPairKick\_involutive},
\texttt{witnessPairKick\_preserves\_fockInner},
\texttt{witnessPairKickLinearEquiv}
 & \NewResult{} \Kernel{} & the normalized Pluecker-phase pair kick is a complex-linear automorphism, exactly reversible, and preserves the finite Fock inner product\\
\texttt{quarticPairTransfer\_isHermitian},
\texttt{pairKick\_eq\_quartic\_add\_offPair}
 & \NewResult{} \Kernel{} & the rank-two quartic transfer is Hermitian and is exactly the kick on the pair sector, with identity action off it\\
\texttt{pairKick\_singleton},
\texttt{witnessPairKick\_two\_particle\_nontrivial}
 & \NewResult{} \Kernel{} & the kick fixes every one-particle basis state but acts nontrivially on the explicit two-particle state\\
\texttt{quarticPairTransfer\_sq\_eq\_project}
 & \NewResult{} \Kernel{} & at unit phase the Hermitian quartic squares exactly to its rank-two pair-sector projector\\
\texttt{bKickL\_CARSupported},
\texttt{heisenStep\_CARSupported},
\texttt{heisenFoldBlocks\_CARSupported}
 & \NewResult{} \Kernel{} & the placed even quartic pair gate is strongly CAR-supported on its declared four-mode set, and every finite unit-phase schedule enlarges support only by the union of acted sets; no metric cone is claimed\\
\texttt{bKickL\_commute\_disjoint},
\texttt{witBlock\_forward\_amplitude},
\texttt{witBlocks\_commute}
 & \NewResult{} \Kernel{} & disjoint placed pair gates commute exactly, while the normalized $3+4\ii$ phase gives a nonzero transfer amplitude inside the acted block\\
\texttt{witness\_conjugate\_restOperators},
\texttt{doubleKick\_return\_amplitude}
 & \NewResult{} \Kernel{} & equal-modulus fields $3+4\ii$ and $5$ have exactly conjugate rest operators, while the two-kick loop returns the pair with amplitude $u_2\overline{u_1}$\\
\texttt{doubleKick\_interference\_amplitude},
\texttt{witness\_survival\_probability}
 & \NewResult{} \Kernel{} & vacuum-referenced interference overlap $(1+u_2\overline{u_1})/2$; exact survival probability $4/5$ at the witnesses versus $1$ for the equal-field control\\
\bottomrule
\end{tabularx}
\end{center}

The controls are part of the result rather than decorative examples.
\begin{center}
\small
\begin{tabularx}{\linewidth}{@{}p{0.22\linewidth}X X@{}}
\toprule
Claim & Positive/nondegenerate control & Failure or boundary control\\
\midrule
Area creates a gap & standard basis pair gives $z=1$ and $B_z^2=\id$
 & collinear pair $(1,2),(3,6)$ gives $z=0$ and $B_z=0$\\
Full null boundary matters & determinant vanishes on every rank-one edge
 & the form $ac$ vanishes on coordinate rays but fails on the diagonal ray\\
Mass enters at turns & a three-step history is nonzero and mass-dependent
 & at zero mass every positive-corner history vanishes\\
Many-step convergence & rational $(k,\mu)=(3/5,4/5)$ satisfies the bound
 & reversing the mass phase converges to the wrong Hamiltonian numerically\\
$3+1$ factor composition & $(1,2,2,3)$ gives $H_4^2=18\id$ and nonzero norm
 & $\alpha_1\alpha_2\neq\alpha_2\alpha_1$ forbids order erasure and naive second order\\
\bottomrule
\end{tabularx}
\end{center}

The principal verification commands are
\begin{verbatim}
lake env lean PhysicsSM/Draft/NullEdge/PluckerMassOperator.lean
lake env lean PhysicsSM/Draft/NullEdge/PluckerMassDynamics.lean
lake env lean PhysicsSM/Draft/NullEdge/CornerWalkEquivalence.lean
lake env lean PhysicsSM/Draft/NullEdge/FermionDoublingAudit.lean
lake env lean PhysicsSM/Draft/NullEdge/ExactCheckerboardPathSum.lean
lake build PhysicsSM.Draft.NullEdge.ExactQuantumWalkDispersion
lake build PhysicsSM.Draft.NullEdge.FixedMomentumManyStepContinuum
lake build PhysicsSM.Draft.NullEdge.BoundedMomentumManyStepContinuum
lake build PhysicsSM.Draft.NullEdge.Compact3Plus1DiracRate
lake env lean PhysicsSM/Draft/NullEdge/Compact3Plus1RefinedWindowRate.lean
lake env lean PhysicsSM/Draft/NullEdge/Local3Plus1RateBridge.lean
lake env lean PhysicsSM/Draft/NullEdge/FiniteWalkPositionConvergence.lean
lake env lean PhysicsSM/Draft/NullEdge/FiniteZModPlancherel.lean
lake env lean PhysicsSM/Draft/NullEdge/Finite3Plus1FourierBridge.lean
lake env lean PhysicsSM/Draft/NullEdge/Finite3Plus1AnalyticSignBridge.lean
lake env lean PhysicsSM/Draft/NullEdge/Pluecker3Plus1ComplexMass.lean
lake env lean PhysicsSM/Draft/NullEdge/ComplexPlueckerLocalWalk.lean
lake env lean PhysicsSM/Draft/NullEdge/ComplexPlueckerRateTransfer.lean
lake env lean PhysicsSM/Draft/NullEdge/FiniteTorus3Plancherel.lean
lake env lean PhysicsSM/Draft/NullEdge/FiniteTorus3WalkWavepacket.lean
lake env lean PhysicsSM/Draft/NullEdge/Finite3Plus1BrillouinAudit.lean
lake env lean PhysicsSM/Draft/NullEdge/StationaryAmplitudeNoGo.lean
lake env lean PhysicsSM/Draft/NullEdge/StationaryAmplitudeLiveAxisNoGo.lean
lake env lean PhysicsSM/Draft/NullEdge/FiniteWalkOnsiteEquivalenceObstruction.lean
lake env lean PhysicsSM/Draft/NullEdge/LocalQCAProperties.lean
lake env lean PhysicsSM/Draft/NullEdge/CountableL2WavepacketConvergence.lean
lake env lean PhysicsSM/Draft/NullEdge/ContinuumL2MultiplierBridge.lean
lake env lean PhysicsSM/Draft/NullEdge/CompactSupportL2WalkBridge.lean
lake env lean PhysicsSM/Draft/NullEdge/VariablePlueckerLocalWalk.lean
lake env lean PhysicsSM/Draft/NullEdge/VariablePlueckerPhaseConnection.lean
lake env lean PhysicsSM/Draft/NullEdge/ComplexPlueckerCheckerboardPathSum.lean
lake env lean PhysicsSM/Draft/NullEdge/GlobalPhaseWindingNoGo.lean
lake env lean PhysicsSM/Draft/NullEdge/PlueckerWindingDerived.lean
lake env lean PhysicsSM/Draft/NullEdge/StrictQCAMinimalArchitecture.lean
lake env lean PhysicsSM/Draft/NullEdge/CommutatorWilsonStrictnessKill.lean
lake env lean PhysicsSM/Draft/NullEdge/TemporalBlockingRG.lean
lake env lean PhysicsSM/Draft/NullEdge/FullBlochSplitPlus.lean
lake env lean PhysicsSM/Draft/NullEdge/FullBlochSplitMinus.lean
lake env lean PhysicsSM/Draft/NullEdge/FullBlochZeroClassification.lean
lake env lean PhysicsSM/Draft/NullEdge/FiniteCARFockBasic.lean
lake env lean PhysicsSM/Draft/NullEdge/FiniteCARSecondQuantization.lean
lake env lean PhysicsSM/Draft/NullEdge/PlueckerQuarticInteraction.lean
lake env lean PhysicsSM/Draft/NullEdge/PlueckerPhaseObservable.lean
lake env lean PhysicsSM/Draft/NullEdge/PlueckerPhaseDefectSpectrum.lean
lake env lean PhysicsSM/Draft/NullEdge/ChangingLatticePDECore.lean
lake env lean PhysicsSM/Draft/NullEdge/Carrier/PluckerScaleSelectionNoGo.lean
lake build PhysicsSM.Draft.NullEdge.OvernightTheoryAxiomGuard
python Scripts/sim/null_edge_regulator_benchmark.py
python Scripts/sim/plot_null_spinor_dirac_bands.py
python Scripts/sim/plot_null_edge_paper_figures.py
python Scripts/publication/verify_null_edge_paper_a.py
\end{verbatim}
The cited modules contain no proof handoff markers or compiled-evaluator proofs.
The working tree is not yet an archival release.  A clean commit, source
archive, theorem ledger, and DOI remain submission requirements.

\subsection*{What mechanization changed}

The formal artifact is scientifically useful only where it improves the
mathematics.  In this development it forced ten corrections that are visible
in the paper.  First, it separated the first-order rest energy $\mu$ from the
quadratic invariant $\mu^2=\det P$ instead of allowing one symbol to drift
between operator orders.  Second, exact unitarity exposed the finite-spacing
corner ratio $\tan(a\mu)$, replacing the tempting but false finite identity
$a\mu=\tan(a\mu)$.  Third, proving an iterated product estimate separated a
one-step Taylor expansion from genuine uniform many-step control.  These are
not the only bookkeeping gains.  Fourth, the exact $3+1$ composition preserved
the noncommuting x/y/z/mass factor order instead of identifying reversed local
 updates.  Fifth, the $1+1$ full-zone audit separated the zero-quasienergy cone
 from its high-quasienergy Floquet partner.  Sixth, the $3+1$ high-symmetry audit
 disproved the tempting global reading of continuum isotropy: the tangent is
 isotropic, three nonzero corners exactly alias the origin, and at body center
 the massive step retains explicit $+1$ and $-1$ modes for every mass angle.
 Seventh, the exact Laurent residual ruled out the most economical proposed
 repair: a nonzero stationary amplitude cannot coexist with origin
 normalization, exact unitarity, and a full involutory tangent inside one
 degree-one nearest-neighbor factor.
 Eighth, the inhomogeneous phase calculation separated a global basis rotation
 from a genuine endpoint connection and proved that a global real lift cannot
 carry nonzero winding.  Ninth, exact temporal blocking showed that the
 mass-only group law closes while the noncommuting spatial split generates a
 larger ordered-word family.  Tenth, the radial derivative of the positive
 Pluecker action proved that its homogeneity collapses rather than selects a
 nonzero scale.
 These are the formalization contributions; the number of declarations is not.

\section{What is forced, derived, realized, and open}

The central scientific distinction is not between bold and cautious language.
It is between different logical verbs.

\begin{center}
\begin{tabularx}{\linewidth}{@{}p{0.16\linewidth}X@{}}
\toprule
Verb & Claim in this paper\\
\midrule
Forced & A quadratic invariant vanishing on the full null boundary is
proportional to the determinant; a nonnegative first-order gap is proportional
to its square root.\\
Derived & The spinor pair gives $z$; $z$ gives $B_z$; $B_z$ gives the Dirac
square, rest gap, unitary turn gate, and orientation-sensitive corner phase;
a site field $z(x)$ gives its own patched link data with an exact
$2\pi$-integer winding.\\
Realized & The split-step null walk exhibits the construction exactly and
converges uniformly on bounded $1+1$ momentum boxes; a finite local $3+1$
 successive-axis walk has an ordered complex-Pluecker symbol with the same
 explicit rate, a proved quartically growing-window consistency theorem,
 exact three-torus Plancherel control, and a classified
 high-symmetry failure set.\\
 Open & A walk-specific changing-lattice Fourier/PDE limit, a replacement
 $3+1$ regulator that removes the corner aliases and body-center modes,
 beyond the excluded degree-one stationary-amplitude ansatz, reduction of
 the derived link data to a corrected non-winding-only defect index and stable
 localized mode, a graph-metric cone for the exact placed pair layer and its composition with the spatial
 free walk, an operator-exponential/action law, and a continuum interaction
 theory.\\
\bottomrule
\end{tabularx}
\end{center}

What can an assigned scalar mass mimic?  Everything that depends on $|z|$
alone: after setting $m=|z|$, the constant-coefficient spectrum, the
dispersion, the velocity bound, and the uniform many-step limit are
identical, and the paper says so without embarrassment.  What it cannot
mimic is now a theorem list rather than a narrative.  It has no input from
which the rest operator is derived rather than declared (the classification
theorem makes ``derived'' precise); it carries one real coin where the exact
history sum carries $z$ and $\bar z$ on opposite orientations; it produces
no endpoint connection, no integer link winding, and no patched transition
data from a varying phase unless that field is separately supplied; and a
scalar mass, being one real number, cannot reproduce the exact
$4/5$-versus-$1$ interference split that separates two equal-modulus complex
Pluecker fields --- matching it requires importing exactly the complex datum
this paper derives, together with the pair interaction that reads it.  At
constant $z$ the distinction is structural: the economy of deriving mass and
interaction phase from one object.  Once the phase varies in space, the
distinction becomes patched link data with an integer winding; whether the
free walk forces the interaction that would make the constant-$z$ phase a
free-theory observable is open.

There is now also a free two-site spectral consequence.  For two equal-modulus
rest blocks coupled by the displayed unit edge phase, the exact $4\times4$
Hamiltonian obeys a scalar quadratic polynomial whose splitting is controlled
by the parallel-transported mismatch
$\Delta=z_R-(\overline w)^2z_L$.  Its lower squared-energy level vanishes
exactly at $t=|z_L|$ and
$z_R=-(\overline w)^2z_L$ (for $t\geq0$), while a common phase on both sites
is removed by an explicit unitary conjugacy.  The equal-modulus hypothesis is
load-bearing: the formal development contains a concrete unequal-modulus
counterexample.  This is a phase-sensitive finite free spectrum, not yet a
topological-protection or localization theorem.

The paper proves a finite free construction and one finite phase-sensitive
quartic extension.  It does not prove that every
massive particle is composite in the ordinary particle-physics sense.  The
``constituents'' here are null momentum factors of a finite spinor description,
not additional asymptotic particles.  Nor does the theorem compute an absolute
mass scale: rescaling the input spinors rescales $z$ and hence $\mu$.  The
quartic-action no-go sharpens this from an omission to a proved boundary for
the natural positive homogeneous action.

The result does support a precise statement: a timelike rest gap can be encoded
entirely by massless rank-one momentum primitives, and the gap is exactly the
area opened between their spinor directions.  In the derived walk, microscopic
translation occurs along null channels while subluminal drift and a rest gap
arise from coherent direction changes.  This is an exact property of the
finite model, not yet a universal ontology of nature.

\section{Interpretation and open problems}

The larger null-information program studies a four-term finite carrier square,
positive cohomology, closure interactions, and discrete frame data.  Those
structures motivated the present theorem but are not required to prove it.
They belong in a companion paper unless they acquire a theorem that depends on
$B_z$.

Three claim-discipline boundaries govern everything in this paper and its
companions.  First, every invariant here is a finite-dimensional algebraic
discriminator on an explicit register; we claim no Fredholm or
finite-volume topological index, and any comparison with established
real-space walk indices is made through their closed-form evaluation on
infinite periodic extensions, never as a finite-volume theorem.  Second,
the pair interaction is supplied as an even Hermitian four-mode kick
consistent with the CAR support and layer schedule; we prove its exact
gate, unitarity, and composition behavior, but we do not derive the
interaction from a microscopic principle, and the same gap is open in the
interacting-automaton literature we compare against.  Third, every result
in this paper is machine-checked in Lean~4 with its axiom footprint
pinned; where companion analyses use exact computer-algebra computations
that are not yet formalized, they are tagged as such there and are never
cited here as theorems.

The next problems are sharply ordered.
\begin{enumerate}[leftmargin=2em]
  \item \textbf{Complete the lattice-to-PDE identification.}
  The finite product DFT, its exact $N^{-3}$ normalization, and the
  walk-specific compact-support momentum-space $L^2$ limit are closed.  Control
  the changing lattice/continuum Fourier spaces and identify the limiting
  multiplier with the position-space Dirac PDE propagator.  The abstract
  bulk/tail split and vanishing of the ultraviolet tail under measurable
  monotone band exhaustion are now closed for arbitrary $L^2$ data, and the
  per-step symbol estimate now carries the refined form
  $2B_4^2t^2/n\cdot\exp(|t|B_4/n)$, which retains the small-step factor
  that the earlier bound discarded; the remaining frontier is the
  balancing of a growing momentum window against the Sobolev tail,
  composed with the concrete scaled sampling/interpolation isometry.
  Equations \eqref{eq:plancherel3}, \eqref{eq:countablel2},
  \eqref{eq:continuousl2multiplier}, and \eqref{eq:walkcompactl2} close the
  finite, abstract countable, and compact-support multiplier steps, not that
  final physical identification.
  \item \textbf{Remove the lattice partners without losing the derivation.}
  The no-go \eqref{eq:stationarynogo} rules out the naive degree-one
  stationary-amplitude factor with a full Dirac tangent.  The Wilson
  construction \eqref{eq:wilsonsquare} now removes the unwanted corners and
  gives a uniform massive gap for a nearest-neighbor Hamiltonian.  The open
  task is to retain that spectral success in a strictly finite-range,
  exactly unitary discrete-time update.  Test the enlarged
  internal cells, longer Laurent range, and tilted-projector stay-put
  constructions of the kind realized by Gupta and Short
  \cite{GuptaShortDoubling} while preserving the
  Pluecker-derived coin and exact path dictionary; by the corollary above,
  any such stay-put route must quantify what replaces the involutory
  unit-speed tangent.  In $3+1$, remove the three
  proven even-parity origin
   aliases and both mass-independent body-center modes while retaining locality,
   unitarity, continuum isotropy, and the complex-mass rate, and eliminate
   the residual non-Dirac states that persist in the Gupta--Short family.
   Equations
   \eqref{eq:cornerparity} and \eqref{eq:bodycentermodes} are the regression
   tests.
  \item \textbf{Complete the tetrahedral comparison.}
  Compute the tetrahedral walk's full Bloch symbol in a common blocked unit
  cell and audit its corners, bands, and locality radius.  Constant onsite
  equivalence is already constrained by the alias-preservation no-go; a
  genuine comparison must specify any blocking or coarse-graining map.
  \item \textbf{Generalize the Pluecker rest operator.}
  For larger null families, seek a canonical odd Hermitian operator from
  $w=(\psi_i\wedge\psi_j)_{i<j}$ whose square is the complete area budget and
  whose phase/moduli classification is decomposition-independent.
  \item \textbf{Reduce the derived winding to an index and a localized mode.}
  The position-dependent unitary update, its inverse and causal cone, and exact
  site-dependent chiral connection are closed, and the derivation gap is now
  closed as well: the link increments of any nowhere-zero Pluecker field are
  the principal arguments of its neighbor ratios, their cycle sum is an exact
  $2\pi$-integer winding, a globally lifted real phase still forces zero, and
  an explicit winding-one field is generated by primitive spinors.  The next
  theorem must reduce this derived winding, or a sign/domain wall through the
  collinear zero set, to a finite index operator and exhibit a localized
  protected mode with a displayed stability class.
  \item \textbf{Complete the many-body and interacting lift.}
  The generic finite CAR algebra, a phase-sensitive four-mode quartic
  pair-transfer observable, and an exact norm-preserving pair kick are now
  closed; the quartic operator is Hermitian and the kick/operator restriction
  identity is exact.  The determinant-minor lift now fixes vacuum, agrees with
  the one-particle matrix, preserves number and parity, obeys exact creation
  covariance, and satisfies $\Gamma(1)=1$ and
  $\Gamma(UV)=\Gamma(U)\Gamma(V)$.  Its ordered-minor adjoint law, lifted
  two-sided inverse, and finite Fock inner-product preservation are also exact.
  Annihilation covariance is now exact, and the placed quartic pair layer has
  exact disjoint-block commutation, finite scheduled CAR-support propagation,
  a sequential finite-range graph-neighborhood cone, and one neighborhood
  expansion per pairwise-disjoint circuit layer.  Compose this circuit-depth
  theorem with the spatial one-particle walk, then derive an
  operator-exponential law and the quartic dynamics from one
  field-valued action whose reductions also produce the null-spinor geometry
  and one-particle evolution.
  \item \textbf{Add a scale-bearing selection mechanism and derive couplings.}
  The generic radial theorem now proves that no positive fixed finite
  homogeneous action of positive natural degree has a nonzero stationary
  scale; degree zero is flat.  Therefore a nonzero solution must display a
  refining RG family, competing homogeneity, constraint, ensemble scale, or
  other dimensionful input and quantify which observed scale remains supplied.
  Interaction strengths and gauge representations remain open.
\end{enumerate}

The finite construction encodes spacetime information algebraically through a
fixed Pauli map and null direction grading.  Emergence of a dynamical continuum
geometry remains conjectural.  We therefore do not claim that the paper starts
without spacetime structure; it starts without primitive timelike motion or a
primitive mass term inside the finite model.

\section{Conclusion}

The result readers should retain is compact.  Null spinors open an exterior
area.  That area is the positive rest gap, its square is the momentum
determinant, and its complex orientation is the off-diagonal part of an odd
Hermitian mass operator.  The collinear null boundary is exactly the
gap-closing locus.

The operator does more than rename a scalar.  It produces the relativistic
Dirac square, exponentiates to the exact unitary turn, and fixes the finite
corner ratio, including the directed $z$/$\bar z$ factors in the exact history
sum.  The associated split-step walk converges uniformly on bounded
$1+1$ momentum boxes at an explicit $O(1/n)$ rate, and its transfer-matrix
history representation has a two-weight recursion that scales to the
checkerboard corner ratio $-\ii\tan(a\mu)$.  The
   finite local $3+1$ extension accepts arbitrary position-dependent
   Pluecker profiles, preserves norm exactly, has an explicit inverse and
   causal cone, and turns local phase variation into endpoint links.  Its
   ordered symbol carries the complex phase at the same uniform compact-box
   rate, the three-torus Plancherel identity fixes the exact $N^{-3}$
   wave-packet normalization, and compactly supported momentum-space data
   converge in $L^2$ at every fixed time.  The site field supplies its own
   patched link data with an exact integer winding, generated by primitive
   spinors, and the retained phase --- invisible to the one-particle rest
   spectrum at constant modulus --- is read out by an exact two-particle interference
   probability inside the supplied pair-kick sector.  These statements form
   one finite
   chain from null geometry to rest energy and propagation.  The same audit proves where that chain
   presently stops: the cubic regulator has three exact zero-quasienergy corner
   aliases, its massive body-center step retains exact $0$ and $\pi$
   quasienergies for every mass angle, and for every intermediate mass angle
   those body-center-type points are provably the only zero- or
   $\pi$-quasienergy crossings.  The degree-one stationary-amplitude
   no-go excludes the simplest attempted repair, and its factorized form proves
   that every onsite mass coin leaves all three even-parity aliases intact.
   A Wilson term removes the
   non-origin corners and opens a uniform massive gap at the local-Hamiltonian
   level; converting that success into a strictly finite-range discrete-time
   regulator remains open.  Exact temporal blocking further shows that the
   mass subgroup closes while the full split family generates a larger ordered
   monoid, and the homogeneous-action no-go proves that a nonzero absolute
   mass needs an additional scale-bearing input.

What remains is equally definite: a changing-lattice position-space PDE limit, a local
 strictly finite-range unitary regulator outside the degree-one
 stationary-amplitude no-go that matches the Wilson Hamiltonian's spectral
 success without sacrificing the derived coin, a blocked
 comparison with the tetrahedral construction, a patched defect/index theorem,
 and a second-quantized interacting extension.  Success would extend a machine-verified finite
rest-gap construction into a broader discrete relativistic dynamics.  Failure
at any named step would locate the boundary of the idea rather than dissolve
it into analogy.

\begin{thebibliography}{99}

\bibitem{ElvangHuang}
H.~Elvang and Y.-t.~Huang,
\emph{Scattering Amplitudes}, arXiv:1308.1697.

\bibitem{ArkaniHamedMassive}
N.~Arkani-Hamed, T.-C.~Huang, and Y.-t.~Huang,
\emph{Scattering Amplitudes For All Masses and Spins},
JHEP 11 (2021) 070, arXiv:1709.04891.

\bibitem{DittmaierWvdW}
S.~Dittmaier,
\emph{Weyl--van der Waerden formalism for helicity amplitudes of massive
particles}, Phys. Rev. D 59 (1998) 016007, arXiv:hep-ph/9805445.

\bibitem{FeynmanHibbs}
R.~P.~Feynman and A.~R.~Hibbs,
\emph{Quantum Mechanics and Path Integrals}, McGraw--Hill, 1965.

\bibitem{JacobsonSchulman}
T.~Jacobson and L.~S.~Schulman,
\emph{Quantum stochastics: the passage from a relativistic to a
non-relativistic path integral}, J. Phys. A 17 (1984) 375.

\bibitem{BisioDarianoTosini}
A.~Bisio, G.~M.~D'Ariano, and A.~Tosini,
\emph{Quantum field as a quantum cellular automaton: the Dirac free evolution
in one dimension}, Ann. Phys. 354 (2015) 244, arXiv:1212.2839.

\bibitem{DarianoPathIntegral}
G.~M.~D'Ariano, N.~Mosco, P.~Perinotti, and A.~Tosini,
\emph{Path-integral solution of the one-dimensional Dirac quantum cellular
automaton}, Phys. Lett. A 378 (2014) 3165--3168, arXiv:1406.1021.

\bibitem{MallickSplitStep}
A.~Mallick and C.~M.~Chandrashekar,
\emph{Dirac cellular automaton from split-step quantum walk},
Sci. Rep. 6 (2016) 25779, arXiv:1509.08851.

\bibitem{SucciQW}
S.~Succi, F.~Fillion-Gourdeau, and S.~Palpacelli,
\emph{Quantum lattice Boltzmann is a quantum walk},
EPJ Quantum Technol. 2 (2015) 12.

\bibitem{FreeQCA}
A.~Bisio, G.~M.~D'Ariano, P.~Perinotti, and A.~Tosini,
\emph{Free quantum field theory from quantum cellular automata: Derivation of
Weyl, Dirac and Maxwell quantum cellular automata},
Found. Phys. 46 (2016) 1032--1081, arXiv:1601.04832.

\bibitem{NzonganiTetra}
U.~Nzongani, N.~Eon, I.~Marquez-Martin, A.~Perez, G.~Di Molfetta, and
P.~Arrighi,
\emph{Dirac quantum walk on tetrahedra},
Phys. Rev. A 110 (2024) 042418, arXiv:2404.09840.

\bibitem{AsbothObuse}
J.~K.~Asboth and H.~Obuse,
\emph{Bulk-boundary correspondence for chiral symmetric quantum walks},
Phys. Rev. B 88 (2013) 121406(R), arXiv:1303.1199.

\bibitem{GNVW}
D.~Gross, V.~Nesme, H.~Vogts, and R.~F.~Werner,
\emph{Index theory of one dimensional quantum walks and cellular automata},
Commun. Math. Phys. 310 (2012) 419--454, arXiv:0910.3675.

\bibitem{GuptaShortDoubling}
C.~Gupta and A.~J.~Short,
\emph{Fermion doubling in Dirac quantum walks},
Phys. Rev. A, accepted (2026), arXiv:2601.15885,
doi:10.1103/kkz3-fthm.

\bibitem{MlodinowBrun}
L.~Mlodinow and T.~A.~Brun,
\emph{Discrete spacetime, quantum walks, and relativistic wave equations},
Phys. Rev. A 97 (2018) 042131, arXiv:1802.03910.

\bibitem{ChildsTrotter}
A.~M.~Childs, Y.~Su, M.~C.~Tran, N.~Wiebe, and S.~Zhu,
\emph{A theory of Trotter error}, Phys. Rev. X 11 (2021) 011020,
arXiv:1912.08854.

\bibitem{Lean4}
L.~de Moura and S.~Ullrich,
\emph{The Lean 4 theorem prover and programming language},
CADE-28, Lecture Notes in Computer Science 12699 (2021) 625--635.

\bibitem{Mathlib}
The Mathlib Community,
\emph{The Lean mathematical library}, Proceedings of CPP 2020, 367--381.

\bibitem{PhysLean}
J.~Tooby-Smith,
\emph{HepLean: Digitalising high energy physics}, arXiv:2405.08863.

\bibitem{DouglasQFTLean}
M.~R.~Douglas, S.~Hoback, A.~Mei, and R.~Nissim,
\emph{Formalization of QFT}, arXiv:2603.15770.

\end{thebibliography}

\end{document}
